% =====================================================================
% MIKU SCI 小论文 — 中文整合初稿 v1
% 整合自 build/sections/{abstract,intro,problem,method,theory,experiment,related,conclusion}.tex
% 数据冻结见 SPEC.md B 节，叙事见 C 节，文风见 D 节，红线见 E 节。
% 自包含可编译骨架，编译命令： latexmk -xelatex v1.tex
% =====================================================================
\documentclass[a4paper,11pt]{ctexart}

\usepackage{amsmath,amssymb,amsthm}
\usepackage{booktabs}
\usepackage{graphicx}
\usepackage{enumitem}
\usepackage{geometry}
\usepackage[ruled,vlined,linesnumbered]{algorithm2e}
\usepackage[backend=biber,style=gb7714-2015,gbpub=false,
            uniquename=init,uniquelist=true,
            maxcitenames=2,mincitenames=1,
            sorting=none,doi=false,url=false,
            natbib=true]{biblatex}

\geometry{margin=2.5cm}

% ---- 定理类环境（中文，按节编号）----
\theoremstyle{definition}
\newtheorem{proposition}{命题}[section]
\newtheorem{lemma}{引理}[section]
\newtheorem{theorem}{定理}[section]
\renewcommand{\proofname}{证明}

% ---- 行距与列表 ----
\linespread{1.15}
\setlist{nosep}
\emergencystretch=3em  % 容纳无法断行的长等宽标识符，避免 overfull

% ---- 参考文献 ----
\addbibresource{references.bib}

\title{显式约束传递通道：弥合多障碍物自动驾驶中的路径--速度信息断裂}
\author{廖嘉旺，孙成娇\thanks{通讯作者。}\\[2pt]\normalsize 湖北文理学院 计算机工程学院}
\date{}

\begin{document}
\maketitle

% PLACEHOLDER_ABSTRACT
\begin{abstract}
\noindent
路径与速度解耦是量产自动驾驶的主流范式，但两阶段间存在结构性信息断裂，路径阶段缺乏时间维度，速度阶段只能被动制动，致使密集场景频繁出现阻塞与过度保守避让。本文提出多障碍物交互运动学统一规划算法（Multi-obstacle Interaction-density Kinematic-prediction Unified planner, MIKU），在解耦接口处建立显式约束传递通道，由三个组件构成。其一以融合碰撞时间、横向重叠率、相对速度、类型与交互密度的多因子威胁度评估替换统一缓冲，输出差异化安全裕度 $\delta_i$；其二以扫描线检测将关联障碍物聚为连通分量，组级求解最大间隙最优通行带，将组合搜索复杂度由 $O(2^k)$ 降至 $O(k\log k)$；其三引入到达时间函数 $\tau(s)$ 与预测轨迹构造时空可通行走廊，经 STDrivableBoundary 注入速度二次规划（Quadratic Programming, QP）。给出纯解耦不可行的充分条件，并以连续性引理与最大间隙最优性定理证明走廊非空时的可行性恢复。MIKU 以 LaneFollowPath 路径边界分支接入 Apollo 11.0 生产级代码，在 8 个仿真场景验证，含 4 个压力场景与 4 个可比场景。压力场景下 Baseline 通过率 0/4，MIKU 4/4，汇总平均速度由 3.53 升至 4.57 m/s，相对提升 29.3\%；可比场景 QP 平均总耗时由 19.16 ms 降至 6.26 ms。

\vspace{0.5em}
\noindent\textbf{关键词：}自动驾驶；路径规划；Apollo；MIKU 算法；组合优化
\end{abstract}
% ====== _parts/intro.tex ======
\section{引言}
\label{sec:intro}

路径--速度解耦（Path-Velocity Decomposition, PVD）是当前量产自动驾驶系统中应用最广泛的轨迹规划范式，其思想自\citet{kant1986pvd}提出以来一直是结构化道路规划的理论起点，即先在Frenet坐标系的空间域规划一条绕开静态障碍物的路径，再在以路程为横轴、时间为纵轴的时空域上规划速度以规避动态障碍物。\citet{fan2018baidu}首次将这一两阶段优化完整部署到工业级的L4自动驾驶栈，使路径层与速度层各自归结为一个低维二次规划（Quadratic Programming, QP）子问题；这一框架被后续的百度Apollo规划器\citep{apollo2024}沿用，Autoware的Lattice规划器、京东配送车ROVER~5.0\citep{wang2022_jfr}等工业系统也采用了类似的分解策略。解耦之所以成为主流，在于它把高维联合优化降解为两个可在凸框架内求解的子问题，从而把规划周期稳定控制在车载实时性所要求的100\,毫秒以内\citep{fan2018baidu}。

解耦带来实时性的同时，也在两个阶段的接口处留下一道结构性的信息断裂。路径阶段建立在静态化假设之上，其构建的可行驶边界不含时间维度，无法表达某个障碍物只在特定时刻占用某段空间这一事实；速度阶段拿到的是一条已经确定的固定路径，面对动态障碍物时只能沿该路径被动地加减速，而无法回头调整路径的横向走向。以一个具体场景为例，自车前方有一动态障碍物，其预测包络在规划时域内横扫过本车道的可通行带。若要在路径阶段显式顾及该障碍物，只能取其在全时域上占用区域的并集作为禁行区，而该并集在某一纵向区间上可能覆盖整条可行带，于是路径子问题在该区间上无可行解；但在任意单一时刻，可行带其实仍有空隙，一次恰当的通行时机选择本可让自车从障碍物经过之前或之后穿过。Apollo的工程实现回避了这一并集投影，转而在路径边界构建中通过静态性判定把动态障碍物整体交给速度阶段处理，路径因而绕开静态障碍物却对动态障碍物视而不见，速度阶段只得制动甚至停车。无论采取哪种处理，路径层无时间维与速度层只能制动这两个表现，都源自同一道信息断裂，并在多障碍物密集场景下被进一步放大\citep{gonzalez2016}。

针对这道断裂，现有研究大致沿两条路线展开，各有其代价。第一条路线放弃解耦，转向时空联合规划。\citet{yang2021_ia_planner}提出IA~Planner，通过瞬时分析将三维时空约束投影为二维空间约束后用模型预测控制（Model Predictive Control, MPC）求解联合轨迹；\citet{hu2023_hybrid_astar_slt}在三维的路程--横向--时间空间中用改进混合A\textsuperscript{*}搜索粗粒度轨迹，再用非线性MPC细化；\citet{qiao2025stjoint}则在结构化道路下引入时间--纵向--横向三维决策并与时空安全走廊的二次规划相结合。这类方法换得了时空一致性，但需重写整条求解链路，计算量随维度成倍增长，难以直接复用现有的生产级框架。第二条路线保留解耦框架，在其内部做局部改进。\citet{han2026_contingency}在路径阶段的QP目标函数中注入障碍物的多模态预测概率作为软代价，以隐式方式传递预测信息，然而软代价只调整目标的权重，并不构成对可行域的约束，无法给出可行性保证；\citet{zhang2023_dp_rcg}并行枚举候选的路径--速度组合并用合作博弈选取最优配对，其枚举复杂度为$O(mn)$，且对原有架构的改动量较大。

本文走第三条路。核心观察是，断裂的根源不在于解耦本身，而在于解耦接口上缺少一条把结构化时序信息从速度维度回传至路径阶段的通道。据此提出多障碍物交互运动学统一规划算法（Multi-obstacle Interaction-density Kinematic-prediction Unified planner, MIKU），在路径与速度的解耦接口处建立一条显式的约束传递通道，以硬约束而非软代价的形式把信息传回路径边界构建阶段，且不改变路径QP与速度QP各自的内部求解结构，计算代价几乎不增加。这条通道由三个相互衔接的组件构成，分别回答三个问题。其一是传哪些，由多因子威胁度评估区分障碍物的威胁程度，映射为差异化的安全裕度$\delta_i$，使通行空间不再被统一裕度一刀切地压缩。其二是空间怎么传，用扫描线分组把同一截面上的障碍物聚为若干组，并在组级上求解最大间隙的最优通行带，配合最优分割的连续性引理与最大间隙最优性定理，把避让方向的组合搜索复杂度由$O(2^k)$降至$O(k\log k)$。其三是时间怎么传，引入到达时间函数$\tau(s)$表示自车行至路程$s$处的预计到达时刻，其值由二阶运动学估计或上一周期的速度曲线给定而非作为优化变量，据此把时变的通行带投影到时空域，构造时空可通行走廊，再经STDrivableBoundary注入速度QP。理论上，本文给出纯解耦路径子问题不可行的一个充分条件，并证明在所构造走廊非空时可行性可以恢复。

围绕这一条通道，本文的工作可归纳如下。在方法层面，提出解耦接口处的显式约束传递通道及其三个组件，统一刻画威胁度过滤、组级最优通行带与时空走廊注入，并保持整体复杂度为$O(n\log n + K)$，其中$K$为路径或时间的离散化规模。在理论层面，明确给出纯解耦失效的充分条件与走廊非空时的可行性恢复保证，并将可行性的讨论收敛在可证明的范围之内。在工程层面，改进沿用Apollo~11.0原有的数据流与QP求解链路，不引入外部求解器，在LaneFollowPath与STDrivableBoundary链路内完成接入，对照基线为Apollo原生的PiecewiseJerkPathOptimizer，下文统称Baseline。在验证层面，于8个仿真场景上对比，其中4个压力场景复现了解耦失效，Baseline通过0/4而MIKU通过4/4，另4个可比场景两者均能通过且轨迹质量相当；汇总来看通过率由Baseline的4/8升至MIKU的8/8，平均速度由3.53\,m/s升至4.57\,m/s，QP平均总耗时由21.56\,ms降至5.56\,ms，并辅以6变体消融与100次蒙特卡洛灵敏度分析。需要说明的是，压力场景刻意设置为暴露解耦失效，相关结论应理解为对失效情形的复现与范围限定，而非常规工况下的普适性结论；本文亦未涉及实车验证，相关局限在第\ref{sec:experiment}节集中讨论。

全文余下部分组织如下。第\ref{sec:problem}节建立问题模型并刻画解耦失效的机理。第\ref{sec:method}节给出MIKU方法及其三个组件。第\ref{sec:theory}节给出理论分析。第\ref{sec:experiment}节报告仿真、消融与灵敏度实验并讨论局限。第\ref{sec:related}节梳理相关工作。第\ref{sec:conclusion}节总结全文。

% ====== _parts/problem.tex ======
% 问题建模与解耦失效分析
% 集成提示：本节使用 proposition 环境，导言区需添加 \newtheorem{proposition}{命题}[section]
% （毕设 preamble 仅定义了 lemma/theorem，整合到 drafts/vN.tex 时补上）。
\section{问题建模与解耦失效分析}
\label{sec:problem}

路径与速度解耦是当前量产规划系统的主流范式，Apollo 的 PublicRoadPlanner、Autoware 以及京东等平台均沿用这一思路。其核心是把三维时空轨迹规划降维为两个顺序求解的二维子问题，先在纵向—横向（station--lateral, SL）空间忽略时间维度规划横向路径 $l(s)$，再在已知路径上沿弧长规划纵向—时间（station--time, ST）空间的速度曲线 $s(t)$。两个阶段以确定性的路径几何为接口，速度阶段拿到固定路径后只能被动调整速度。这一接口截断了时间信息，构成本文关注的结构性信息断裂，也是后续显式约束传递通道所要修复的对象。

\subsection{解耦两子问题的形式化}
\label{subsec:decouple_formalize}

在 Frenet 坐标系下，第 $i$ 个障碍物在 SL 空间中的投影记为轴对齐矩形
\begin{equation}
    O_i = [s_i^{-},\, s_i^{+}] \times [l_i^{-},\, l_i^{+}],
    \label{eq:prob_obs_proj}
\end{equation}
其中横向坐标以参考线为原点、左正右负。为保证安全，在障碍物边界外预留综合安全裕度 $\delta = W_{ego}/2 + d_{buf}$，其中 $W_{ego}$ 为自车宽度，$d_{buf}$ 对应 Apollo 中 \texttt{obstacle\_lat\_buffer} 参数，默认取 $0.4\,\mathrm{m}$。由此定义两侧安全边界
\begin{equation}
    u_i = l_i^{-} - \delta, \qquad v_i = l_i^{+} + \delta,
    \label{eq:prob_uv}
\end{equation}
$u_i$ 为从右侧通过时自车参考点横向坐标的上界，$v_i$ 为从左侧通过时的下界。为每个在截面 $s$ 处活跃的障碍物分配二值避让方向决策 $d_i \in \{L, R\}$，记右绕与左绕集合分别为 $\mathcal{R} = \{i \mid d_i = R\}$ 与 $\mathcal{L} = \{i \mid d_i = L\}$，则该截面的路径横向边界与通行带宽度为
\begin{align}
    l^{+}(\mathbf{d}) &= \min\Big(l_{road}^{+},\; \min_{i \in \mathcal{R}} u_i\Big), \label{eq:prob_upper} \\
    l^{-}(\mathbf{d}) &= \max\Big(l_{road}^{-},\; \max_{i \in \mathcal{L}} v_i\Big), \label{eq:prob_lower} \\
    W(\mathbf{d}) &= l^{+}(\mathbf{d}) - l^{-}(\mathbf{d}). \label{eq:prob_band}
\end{align}
当 $W(\mathbf{d}) > 0$ 时存在横向可通行空间，$W(\mathbf{d}) \le 0$ 时路径被阻塞。

真实场景中障碍物随时间运动。引入自车到达纵向位置 $s$ 的预估时刻 $\tau(s)$，障碍物投影随之成为时间的函数 $O_i(\tau(s))$，安全边界相应写为 $u_i(s) = l_i^{-}(\tau(s)) - \delta$ 与 $v_i(s) = l_i^{+}(\tau(s)) + \delta$。需要强调，$\tau(s)$ 由速度曲线反推或二阶运动学模型估计得到，在路径层作为已知输入而非优化决策变量。理想的联合路径决策目标是在给定 $\tau(s)$ 下使全路径最小通行带宽度最大化
\begin{equation}
    \mathbf{d}^{*} = \arg\max_{\mathbf{d}} \;\min_{s} \big[\,l^{+}(s, \mathbf{d}, \tau) - l^{-}(s, \mathbf{d}, \tau)\,\big],
    \label{eq:prob_joint}
\end{equation}
其决策变量仅为方向向量 $\mathbf{d} = (d_1, \ldots, d_n)$，搜索空间为 $2^n$，暴力枚举复杂度 $O(2^n)$ 随障碍物数量指数增长，在 Apollo $100\,\mathrm{ms}$ 规划周期约束下无法满足实时性要求。

解耦范式并不求解式~\eqref{eq:prob_joint} 的全局形式，而是把它拆成两个顺序子问题。路径子问题在 SL 空间求横向边界与决策 $\mathbf{d}$，由 \texttt{PathBoundsDecider}（Apollo~11.0 中由 \texttt{PathBoundsDeciderUtil} 承载）构建通行带，再由 \texttt{PiecewiseJerkPathOptimizer} 在带内求最优偏移曲线 $l(s)$；速度子问题接收固定路径 $l(s)$，由 \texttt{SpeedBoundsDecider} 将障碍物映射到 ST 图构建约束，再由 \texttt{PiecewiseJerkSpeedOptimizer} 以二次规划（Quadratic Programming, QP）求速度曲线 $s(t)$。两者通过路径几何单向串联，路径阶段输出的几何形态被传递给速度阶段，而路径决策所依赖的时间有效性前提，例如某条路径以来某动态障碍物在特定时刻已移开，并未随之传递。

\subsection{全时刻并集投影丢失时间维度}
\label{subsec:union_projection}

路径子问题忽略时间，其障碍物足迹只能取某个与时间无关的占据区域。在不破坏避障安全性的前提下，唯一稳妥的选择是全时刻并集投影
\begin{equation}
    \bar{O}_i = \bigcup_{t \in [0,\,T]} O_i(t),
    \label{eq:prob_union}
\end{equation}
即把障碍物在整个规划时域内扫过的位置全部并入一个静止足迹。式~\eqref{eq:prob_union} 是最保守的时间无关替代，它抹掉了障碍物随时间移动的全部信息。Baseline 在工程上采用的是这一思路的退化形式，\texttt{PathBoundsDecider} 对静态障碍物仅取 $t{=}0$ 时刻的空间快照，动态障碍物则通过 \texttt{IsWithinPathDeciderScopeObstacle} 中的 \texttt{IsStatic()} 过滤逻辑直接排除于路径决策之外，转交速度阶段在 ST 图中处理。Apollo 开发团队在该过滤函数源码中留有 \texttt{TODO} 注释，提示部分当前低速或静止的障碍物可能在短时间内重新运动，说明该过滤逻辑附近存在未完成的工程考虑。无论取并集足迹 $\bar{O}_i$ 还是 $t{=}0$ 快照，路径层都无从感知障碍物的预测轨迹。其直接后果是，横穿行人在 $t{=}0$ 占据车道中心时路径层无法利用其后续移动，速度层被迫做出不必要的减速乃至停车，而自车实际到达时行人往往已横穿至车道一侧，间隙早已打开，原本无需停车。

\subsection{解耦可行域是联合可行域的子集}
\label{subsec:feasibility_subset}

把联合规划在三维时空中的可行轨迹集合记为 $\mathcal{F}_{joint}$，它由含时间维度的约束界定，自车的横向选择可以随时变约束协同调整。解耦规划分两步进行，路径子问题先在不含时间的约束下确定横向边界，速度子问题再在固定路径上调整速度，两层互不回溯。于是解耦可达的轨迹集合 $\mathcal{F}_{dec}$ 等于路径可行域与速度可行域的笛卡尔积，其路径分量受式~\eqref{eq:prob_union} 的时间无关足迹约束。由于任一解耦可行轨迹同时满足联合问题的约束，而联合问题允许的时变协同在解耦中被禁止，二者满足包含关系
\begin{equation}
    \mathcal{F}_{dec} \subseteq \mathcal{F}_{joint}.
    \label{eq:prob_subset}
\end{equation}
当路径层以并集投影或 $t{=}0$ 静态快照取代时变投影时，路径可行域被进一步收缩，包含关系可退化为真子集 $\mathcal{F}_{dec} \subsetneq \mathcal{F}_{joint}$，由此出现解耦不可行而联合可行的情形。横穿行人即为一例，静态投影下路径层判定 BLOCKED，而沿 $\tau(s)$ 利用时变投影后中心间隙打开，路径重新可行。

\subsection{纯解耦不可行性的充分条件}
\label{subsec:infeasibility_condition}

解耦失效来自两个相互独立的机制，一是空间维度上逐障碍物贪心选向，二是时间维度上的静态投影。Baseline 在 \texttt{PathBoundsDecider} 中逐障碍物按其中心与参考值的相对位置独立判定绕行方向，决策一旦在循环中做出便不可回退，更新后若某截面 $l^{+}(s) < l^{-}(s)$ 即判定 BLOCKED，调用 \texttt{TrimPathBounds} 截断路径并触发 FallbackPath 停车。该贪心策略在 $2^k$ 种方向组合中仅检查其中一种，$k$ 为该截面活跃障碍物数。这两个机制各自给出纯解耦输出 BLOCKED 的一个充分条件，其形式化陈述与证明在第\ref{sec:theory}节给出，此处只作机理说明。

时间维度上，当动态障碍物的全时刻并集投影在某纵向区间上覆盖了去除半车宽后的整段可行区间时，时间盲的路径子问题在该区间无宽度足够的通行带，而沿到达时间 $\tau(s)$ 的时变投影下通行带仍可能为正，该充分条件由命题\ref{prop:decoupled_infeasible}刻画。空间维度上，逐障碍物贪心可能选中分裂式组合而遗漏可行解，而遍历全部 $k+1$ 个候选分割的最大间隙策略不会输出此类虚假 BLOCKED，这一点由定理\ref{thm:max_gap}保证。两个条件均为充分而非必要，贪心选向在多数场景与全局最优一致，仅在特定对称或密集配置下选中不可行组合，对两障碍物的四种方向组合而言只有 RL、LR 两种可能不可行。当且仅当全部 $2^k$ 种方向组合且全部到达时刻均给出非正间隙时，BLOCKED 才对应真实的联合不可行。换言之，Baseline 输出 BLOCKED 是不可行性的充分条件而非必要条件，其失效本质是解耦可行域 $\mathcal{F}_{dec}$ 相对 $\mathcal{F}_{joint}$ 的真子集收缩，既来自空间维度遗漏的可行组合，也来自时间维度抹除的时变窗口。

上述分析表明，Baseline 的次优性沿安全裕度、空间协调与时间投影三个维度展开，其共同根源在于解耦接口处的信息断裂，即统一安全裕度、逐障碍物独立决策与全时刻并集投影使路径层既缺乏分级机制又缺乏全局与时间视角。MIKU 不另起联合规划，而是在这一接口处建立显式约束传递通道，由表~\ref{tab:prob_channel} 所列三个组件分别回答传哪些信息、空间约束怎么传、时间约束怎么传，方法构造与理论保证在后续各节展开。

\begin{table}[htbp]
\centering
\caption{解耦失效的三个维度与 MIKU 显式约束传递通道的对应}
\label{tab:prob_channel}
\begin{tabular}{p{2.4cm} p{4.6cm} p{4.6cm}}
\hline
失效维度 & 解耦接口处的断裂 & 通道组件的回答 \\
\hline
安全裕度 & 所有障碍物共用统一 $\delta$，低威胁障碍物压缩通行带 & 传哪些，多因子威胁度过滤给出差异化裕度 $\delta_i$ \\
空间协调 & 逐障碍物独立贪心选向，遗漏可行的全局组合 & 空间怎么传，扫描线分组与组级最大间隙通行带 \\
时间投影 & 全时刻并集或 $t{=}0$ 快照抹除预测轨迹 & 时间怎么传，$\tau(s)$ 时变投影构造时空可通行走廊注入速度 QP \\
\hline
\end{tabular}
\end{table}

% ====== _parts/method.tex ======
% method.tex — MIKU 方法部分
% 依赖整合稿 preamble 提供：amsmath、algorithm2e（\KwIn/\KwOut/\KwRet）、
% theorem 环境 lemma/theorem（中文“引理”“定理”）、ctex。
% 本片段不重复定义环境，仅使用。

\section{MIKU 方法与显式约束传递通道}
\label{sec:method}

第\ref{sec:problem}节已经表明，Apollo 路径--速度解耦架构的瓶颈不在于某个模块的参数整定，而在于两个阶段之间的接口本身丢弃了时间维度，路径阶段只看 $t=0$ 的静态快照，速度阶段拿到一条已经固定的路径后只能被动制动。多障碍物交互运动学统一规划算法（Multi-obstacle Interaction-density Kinematic-prediction Unified planner, MIKU）不重写这一架构，而是在解耦接口处建立一条显式的约束传递通道，让路径阶段已经掌握的安全与时序信息能够完整地抵达速度阶段。

这条通道需要回答三个问题，本节的三个组件分别对应其中之一。其一是传递什么，即把 Apollo 一刀切的安全裕度替换为按威胁程度分级的差异化裕度 $\delta_i$，使真正危险的障碍物得到不低于原方案的保护，低威胁障碍物则归还被无谓占用的横向空间。其二是横向信息如何传递，即在单个纵向截面上把多障碍物的避让决策从逐个独立的贪心改为组级协调，以一次线性扫描求出全局最优的通行带。其三是时间信息如何传递，即以到达时间函数 $\tau(s)$ 为桥梁，把动态障碍物的预测轨迹投影到自车实际到达时刻的横向位置，构造时空可通行走廊，并经 \texttt{STDrivableBoundary} 把走廊约束注入速度二次规划（Quadratic Programming, QP）。三个组件不是三项并列的贡献，而是同一条通道上前后衔接的三段，它们共享统一的安全边界记号 $u_i$、$v_i$ 并复用同一套间隙最优性结论。

\subsection{多因子威胁度与差异化安全裕度}
\label{sec:method_threat}

Apollo 的 \texttt{PathBoundsDecider} 在 \texttt{GetBufferBetweenADCCenterAndEdge} 中对所有障碍物返回同一套安全裕度，即半车宽再加固定的额外缓冲，不区分障碍物的类型、速度与横向偏移。其上仅有的分级来自 \texttt{nudge\_l\_buffer} 以 $0.5$\,m/s 速度阈值所做的静态与非静态二分，静态障碍物取 $0.3$\,m、非静态取 $0.4$\,m。这一粒度既无法区分轨迹突变的行人与可预测的同向车辆，也无法区分正面逼近与横向远离的障碍物。在窄路场景中，统一缓冲会把两侧净间距同时压到小于车宽，使 \texttt{PathBoundsDecider} 输出零宽通行带而判定 BLOCKED。

MIKU 以一个归一化的威胁度 $\Theta_i$ 取代上述二分，对每个障碍物综合五个因子加权得到

\begin{equation}
    \Theta_i = w_1 f_{TTC}(i) + w_2 f_{overlap}(i) + w_3 f_{vel}(i) + w_4 f_{type}(i) + w_5 f_{inter}(i),
    \label{eq:threat_score}
\end{equation}

其中各因子均归一化至 $[0,1]$，权重之和为一，故 $\Theta_i \in [0,1]$。权重取 $w_1=0.30$、$w_5=0.25$、$w_2=0.20$、$w_3=0.15$、$w_4=0.10$，呈 $w_1>w_5>w_2>w_3>w_4$ 的次序。该次序先依据各因子对多障碍物场景风险的区分度做定性排序，再在对称路锥与横穿行人等典型场景上核对威胁度排序是否符合驾驶直觉后微调确定。碰撞时间与交互密度对场景风险的区分度最高，赋予较大权重；障碍物类型在同一场景内变化较小，赋予较小权重。

五个因子各刻画一种风险来源。碰撞时间（Time To Collision, TTC）因子衡量纵向碰撞的紧迫性，

\begin{equation}
    \mathrm{TTC}_i =
    \begin{cases}
        \dfrac{s_i - s_{ego}}{\dot{s}_{ego} - \dot{s}_i}, & s_i > s_{ego} \ \text{且}\ \dot{s}_{ego} > \dot{s}_i, \\[0.6em]
        +\infty, & \text{其它},
    \end{cases}
    \label{eq:ttc}
\end{equation}

障碍物位于自车后方、自车不会追上或相对速度为零等退化情形均令 $\mathrm{TTC}_i=+\infty$，不构成纵向威胁。再以分段线性映射

\begin{equation}
    f_{TTC}(i) =
    \begin{cases}
        1, & \mathrm{TTC}_i \leq T_{crit}, \\[0.3em]
        \dfrac{T_{max} - \mathrm{TTC}_i}{T_{max} - T_{crit}}, & T_{crit} < \mathrm{TTC}_i < T_{max}, \\[0.6em]
        0, & \mathrm{TTC}_i \geq T_{max},
    \end{cases}
    \label{eq:f_ttc}
\end{equation}

把碰撞时间转为威胁度，临界值取 $T_{crit}=2.0$\,s，依据 Minderhoud 与 Bovy 对碰撞时间分布的分析，碰撞时间低于该值时驾驶员已难以靠转向规避；上限取 $T_{max}=7.0$\,s，与 Apollo ST 图的时间视野一致。

其余四个因子分别从横向、速度、类型与密集度补全风险画像。横向重叠率因子 $f_{overlap}(i)=\max\!\big(0,\ \min(l_i^{+}, l_{ego}^{+}) - \max(l_i^{-}, l_{ego}^{-})\big)/W_{ego}$ 度量障碍物与自车走廊在横向上的重叠比例，零表示完全错开、一表示正对自车。相对速度因子 $f_{vel}(i)=\sigma(v_{rel,i}/v_{max})$ 采用 $\sigma(x)=1/(1+e^{-5x})$ 的 Sigmoid 映射，$v_{rel,i}=\dot{s}_{ego}-\dot{s}_i$，选择 Sigmoid 而非线性是因为相对速度过零附近的区分更重要，而极高相对速度之间无需精细区分。类型因子 $f_{type}(i)$ 取离散值，综合运动不确定性与事故后果，行人、骑行者等弱势道路使用者（Vulnerable Road User, VRU）取 $1.0$，机动车取 $0.7$，未知可动目标取 $0.5$，静态障碍物取 $0.3$，交通锥等小尺寸标识物取 $0.15$，其中 VRU 在同等碰撞能量下受伤与致死风险远高于车内乘员，因而施加额外保护。交互密度因子

\begin{equation}
    f_{inter}(i) = \frac{1}{N-1} \sum_{j \neq i} \mathbf{1}\!\left[d_{ij} < d_{cluster}\right] \cdot \frac{d_{cluster} - d_{ij}}{d_{cluster}},
    \label{eq:f_inter}
\end{equation}

以 SL 平面欧氏距离 $d_{ij}$ 与聚类半径 $d_{cluster}=10$\,m 加权统计邻近障碍物的聚集程度，孤立障碍物 $f_{inter}\approx 0$，密集车群中接近 $1$，$N=1$ 时令 $f_{inter}=0$。交互密度高的障碍物通常落入同一连通分量，将在第\ref{sec:method_band}节的纵向重叠分组中被统一协调。

威胁度确定后，把额外缓冲在区间 $[d_{buf,min}, d_{buf,max}]$ 上随威胁度线性插值，再加半车宽合成最终安全裕度

\begin{equation}
    d_{buf,i} = d_{buf,min} + (d_{buf,max} - d_{buf,min})\,\Theta_i, \qquad
    \delta_i = \frac{W_{ego}}{2} + d_{buf,i},
    \label{eq:dynamic_delta}
\end{equation}

额外缓冲区间取 $[0.10, 0.40]$\,m。式\eqref{eq:dynamic_delta}在两个端点上保持安全语义。$\Theta_i \to 1$ 时额外缓冲达到 $0.40$\,m，总裕度与 Apollo 原方案相同或更大，安全性不下降；$\Theta_i \to 0$ 时额外缓冲降至 $0.10$\,m，总裕度仍含完整的 $W_{ego}/2$，仅把原本被无谓占用的横向空间归还给通行带。当所有 $\Theta_i$ 取同一值时，式\eqref{eq:dynamic_delta}退化为统一缓冲方案。差异化裕度随即改写每个障碍物的横向安全边界

\begin{equation}
    u_i = l_i^{-} - \delta_i, \qquad v_i = l_i^{+} + \delta_i,
    \label{eq:dynamic_uv}
\end{equation}

其中 $u_i$ 是自车从右侧通过该障碍物时参考点横向坐标的上界，$v_i$ 是从左侧通过时的下界。差异化的 $\delta_i$ 只改变 $u_i$、$v_i$ 的取值，不破坏第\ref{sec:method_band}节引理与定理对 $u$ 值排序性的依赖，下游的组级最优性结论原样成立。

\subsection{扫描线分组与组级最大间隙通行带}
\label{sec:method_band}

确定了每个障碍物的安全边界之后，第二个问题是如何在横向上把这些约束传递成一条全局协调的通行带。Apollo 的 \texttt{PathBoundsDecider} 在 \texttt{UpdatePathBoundaryBySLPolygon} 中逐障碍物决定 LEFT\_NUDGE 或 RIGHT\_NUDGE，不考虑同一 $s$ 截面上其余活跃障碍物。两侧同时存在障碍物时，左侧障碍物下压上边界、右侧障碍物上抬下边界，容易出现 $l^{+}(s)<l^{-}(s)$，路径边界判 BLOCKED 后由 \texttt{TrimPathBounds} 截断并触发 \texttt{FallbackPath} 停车。$k$ 个障碍物同时活跃时避让方向共有 $2^k$ 种组合，贪心仅检查其中一种，可能恰好选中不可行组合而错过可行解，这种失效集中出现在对称或密集配置。MIKU 把单截面问题整理为按安全边界排序后的间隙选择问题，并先按纵向重叠把障碍物分组，使彼此无关的障碍物分别独立求解。

两个障碍物 $O_i$ 与 $O_j$ 称为纵向重叠，当且仅当各自纵向占据区间满足 $[s_i^{-}, s_i^{+}] \cap [s_j^{-}, s_j^{+}] \neq \varnothing$。该关系对称但不传递。取其传递闭包，把 $n$ 个障碍物划为 $m$ 个连通分量 $C_1, \ldots, C_m$，不同分量的纵向区间并集互不相交，任意 $s$ 截面至多受一个分量影响，分量间的避让决策可独立求解。分组用一次扫描线完成，按 $s_i^{-}$ 升序排列后维护当前分量已见的最大 $s^{+}$ 记为 $s_{max}$，若下一障碍物的 $s^{-}$ 不超过 $s_{max}$ 则并入当前分量并更新 $s_{max}$，否则另起新分量。排序为 $O(n\log n)$，扫描本身 $O(n)$，分组总体 $O(n\log n)$。

连通分量不保证存在所有成员共同活跃的公共截面，链式情形中 $O_1$ 与 $O_2$ 重叠、$O_2$ 与 $O_3$ 重叠，而 $O_1$ 与 $O_3$ 可能完全错开。对链式分量取保守包络处理，把同一分量内所有障碍物视为同时施加约束，在其覆盖区间上取所有 $u_i$ 的最小值与所有 $v_i$ 的最大值。由于任一真实截面的活跃集合都是分量的子集，包络宽度是真实通行带宽度的下界，包络宽度不小于车宽即保证该分量内所有截面均可通行；包络不可行只说明充分条件不成立，并不意味着真实场景必然阻塞。这一处理保证安全，但不声称对一般链式分量给出全局最优。全路径的通行带宽度由最窄分量决定，

\begin{equation}
    \hat{W} = \min_{j=1,\ldots,m} \hat{W}(C_j).
    \label{eq:global_width}
\end{equation}

在单个分量内部，单截面最优分割具有可一次扫描求解的结构，本节给出其构造，形式化陈述与证明集中在第\ref{sec:theory}节。第一步压缩搜索空间。把分量内 $k$ 个障碍物按 $u_i$ 值升序排列后，最优避让方向必呈连续分割，即存在分割点 $p\in\{0,1,\ldots,k\}$ 使排序后前 $p$ 个左绕、后 $k-p$ 个右绕，由此 $2^k$ 个任意方向组合被压缩为 $k+1$ 个候选分割，这一连续性由引理\ref{lem:continuous_partition}保证。其中 $p=0$ 表示全部右绕、自车走最右侧通道，$p=k$ 表示全部左绕、自车走最左侧通道，$1\leq p\leq k-1$ 表示自车从第 $p$ 与第 $p+1$ 个障碍物之间的间隙通过。存在 $u_i$ 相同的障碍物时以 $v_i$ 降序作二级排序键保证唯一性。

第二步把每个候选分割的通行带宽度归约为一个几何间隙。仅按 $u_i$ 排序不保证 $v_i$ 单调，故引入前缀左安全边界 $V_p=\max_{1\leq q\leq p} v_{(q)}$，并对 $p=0,1,\ldots,k$ 定义 $k+1$ 个间隙 $g_p$，其中 $g_0$ 是道路右边界到最右侧障碍物的间隙，$g_k$ 是前 $k$ 个障碍物形成的最左安全边界到道路左边界的间隙，中间的 $g_p$ 是前 $p$ 个障碍物的最左安全边界 $V_p$ 与后续障碍物最右安全边界 $u_{(p+1)}$ 之间的间隙，其精确表达见式\eqref{eq:gap_def}。障碍物横向边界不交错时 $V_p=v_{(p)}$，该定义退化为相邻障碍物之间的普通几何间隙。在所有障碍物的安全裕度扩展未越出道路边界的前提下，分割点 $p$ 对应的通行带宽度恰等于 $g_p$，最优通行带宽度为 $W^{*}=\max_{p} g_p$，即取最大间隙对应的分割点为最优解，这一最大间隙最优性由定理\ref{thm:max_gap}给出。当某障碍物膨胀后越出道路边界时，先按引理\ref{lem:boundary_trim}的越界预处理规则强制其避让方向并移出候选列表，使定理前提成立。

几个边界情形可被算法自然吸收。其一，当 $V_p>u_{(p+1)}$ 时 $g_p<0$，该通道不可通行，扫描比较时被自动跳过。其二，当全部 $k+1$ 个间隙都不足车宽，即 $\max_p g_p < W_{ego}$，说明该截面无可通行路径，这是真实 BLOCKED 而非算法局限。所有 $u_i$、$v_i$ 均按差异化裕度 $\delta_i$ 计算，引理与定理仅依赖 $u$ 值排序性而不依赖 $\delta_i$ 是否为常数，故第\ref{sec:method_threat}节的差异化裕度与本节的最优性结论彼此兼容。单截面经排序加一次扫描求解，复杂度由暴力枚举的 $O(2^k)$ 降为 $O(k\log k)$。

\subsection{到达时间函数与时空可通行走廊}
\label{sec:method_corridor}

前两个组件都在 $t=0$ 的静态截面上工作。第三个组件回答时间信息如何传递，把横向通行性沿时间轴推广到动态障碍物，并跨过解耦接口抵达速度规划。其桥梁是到达时间函数 $\tau(s)$，给出自车抵达纵向位置 $s$ 的预估时刻。强调 $\tau(s)$ 在路径阶段是给定参数而非优化变量，由二阶运动学模型或上周期速度曲线估计，路径阶段的决策变量仍只是方向组合。以匀加速运动学估计时

\begin{equation}
    \tau(s) = \frac{-v_{ego} + \sqrt{v_{ego}^{2} + 2\,a_{ego}\,(s - s_{ego})}}{a_{ego}},
    \label{eq:arrival_time}
\end{equation}

其中 $s_{ego}$、$v_{ego}$、$a_{ego}$ 取自 \texttt{VehicleState}。$a_{ego}\to 0$ 时式\eqref{eq:arrival_time}经 Taylor 展开退化为匀速估计 $\tau(s)\approx (s-s_{ego})/v_{ego}$；$v_{ego}\approx 0$ 且 $a_{ego}\approx 0$ 时令 $\tau(s)=+\infty$，算法退回静态版本。实现上优先利用上周期速度曲线 $s_{prev}(t)$，若其单调不减则以一维查表加线性插值求首次满足 $s_{prev}(t)\geq s$ 的时刻，仅在无历史曲线、曲线非单调或查询越界时回退到式\eqref{eq:arrival_time}。单规划周期内的位置预测误差量级约为 $\tfrac{1}{2}a_{max}\Delta t^{2}$，远小于安全裕度区间 $[0.10, 0.40]$\,m，不对通行带边界判定构成实质影响。

有了 $\tau(s)$，每个动态障碍物的横向占据便取自其在自车到达时刻的预测位置。对障碍物 $i$，在 Prediction 模块给出的未来 $5$\,s 预测轨迹中查询时刻 $\tau(s)$ 的位置并线性插值，投影至参考线 Frenet 坐标得 $(s_i(\tau), l_i(\tau))$，结合几何尺寸得到时变 SL 边界 $[l_i^{-}(\tau(s)), l_i^{+}(\tau(s))]$；若该时刻障碍物纵向位置与 $s$ 的距离超过其纵向尺寸一半，则标记其在 $s$ 处不活跃。预测视野之外保守取最后预测位置不变。当前实现取概率最高的一条轨迹，扩展到多轨迹并集可降低漏判风险但查询次数线性增加。第\ref{sec:method_threat}节的横向重叠率因子同样改用该时变投影，即以 $f_{overlap}(i, \tau(s))$ 评估动态障碍物，静态障碍物仍取 $t=0$ 观测位置。

把第\ref{sec:method_band}节的静态间隙优化逐截面推广到动态情形。在每个 $s$ 截面按 $u_i(\tau(s))$ 排序后定义时变间隙

\begin{equation}
    g_p(s) =
    \begin{cases}
        u_{(1)}(\tau(s)) - l_{road}^{-}(s), & p=0, \\[0.3em]
        u_{(p+1)}(\tau(s)) - V_p(s), & p=1,\ldots,k-1, \\[0.3em]
        l_{road}^{+}(s) - V_k(s), & p=k,
    \end{cases}
    \label{eq:time_varying_gap}
\end{equation}

前缀左安全边界相应改为 $V_p(s)=\max_{1\leq q\leq p} v_{(q)}(\tau(s))$。引理\ref{lem:continuous_partition}与定理\ref{thm:max_gap}在每个 $s$ 截面逐截面成立。要求沿一段 $s$ 固定同一组级分割，全路径最优通行带宽度取

\begin{equation}
    W^{*} = \max_{p}\ \min_{s}\ g_p(s),
    \label{eq:time_varying_opt}
\end{equation}

实现上仍以保守包络处理链式连通分量，保证安全可行但不声称全局最优。选定最优分割 $p^{*}$ 后，把 SL 通行性沿 $\tau(s)$ 抬升至 ST 平面，得到时空可通行走廊

\begin{equation}
    \Omega = \left\{ (s,t) \ \middle|\ g_{p^{*}}(s) \geq W_{ego}\ \text{且}\ t \geq \tau(s) \right\}.
    \label{eq:st_corridor}
\end{equation}

走廊有两种几何特征。其一是禁行切片，若某 $s_k$ 处最优间隙不足车宽，整条纵向列 $\{s_k\}\times\mathbb{R}_{+}$ 被排除，在 $s$ 方向表现为水平缺口。其二是 $\tau(s)$ 左下沿，要求自车不早于 $\tau(s)$ 到达 $s$，为动态障碍物移开预留时间。两者组合使 $\Omega$ 在几何上是非凸集合。需要说明的是，方法的凸性来源不在 $\Omega$ 本身，而在其沿离散时间步的逐时间步线性投影 $s_j^{ub}$，这与 Baseline 在 \texttt{PathTimeHeuristicOptimizer} 给出 YIELD/OVERTAKE 组合后所得的线性盒约束在 QP 结构上一致，因而可在不改动速度求解器内部数据结构的前提下嵌入。

\subsection{经 STDrivableBoundary 注入速度二次规划}
\label{sec:method_inject}

走廊只有进入速度求解器才能影响实际轨迹。走廊的有效性依赖自车实际到达时间不早于 $\tau(s)$，即对所有 $s$ 有 $t_{arrive}(s)\geq\tau(s)$，等价地写成速度约束 $s(t)\leq s_k$ 对所有 $t<\tau(s_k)$ 成立，含义是自车不得在 $\tau(s_k)$ 之前到达 $s_k$。该约束以收紧 QP 上界的方式注入。Apollo 的 \texttt{PiecewiseJerkSpeedOptimizer} 在等时间间隔 $\Delta t = 0.1$\,s 上以纵向位置 $s_j$、速度 $\dot{s}_j$、加速度 $\ddot{s}_j$ 为决策变量，目标函数惩罚跟踪误差与加加速度，约束含运动学等式约束、物理限制以及 \texttt{SpeedBoundsDecider} 给出的 ST 边界 $s_j^{lb,\mathrm{SBD}} \leq s_j \leq s_j^{ub,\mathrm{SBD}}$，整体写成

\begin{equation}
    \min_{\mathbf{x}}\ \tfrac{1}{2}\mathbf{x}^{\top}\mathbf{H}\mathbf{x} + \mathbf{f}^{\top}\mathbf{x}
    \quad \text{s.t.}\quad \mathbf{A}_{eq}\mathbf{x} = \mathbf{b}_{eq},\ \ \mathbf{x}^{lb} \leq \mathbf{x} \leq \mathbf{x}^{ub},
    \label{eq:qp_compact}
\end{equation}

由 OSQP 求解。每个动态障碍物取其 SL 投影首次进入候选走廊的纵向位置 $s_i^{0}$ 作为关键位置，连同对应到达时刻组成时间清单 $\mathcal{T} = \{(s_k, \tau_k)\}$，其规模等于动态障碍物数 $|\mathcal{D}|$，远小于时间步数 $K$。走廊约束随即对上界做一次收紧

\begin{equation}
    s_j^{ub} \leftarrow \min\!\left(s_j^{ub},\ s_k\right), \quad \forall\, j: t_j < \tau_k,\ \ \forall\, (s_k, \tau_k) \in \mathcal{T}.
    \label{eq:tighten_ub}
\end{equation}

该操作只改写 $\mathbf{x}^{ub}$ 中对应分量的值，保持 $\mathbf{H}$、$\mathbf{A}_{eq}$ 与稀疏模式完全不变，对纵向占据范围大的障碍物可两端采样把 $\mathcal{T}$ 扩为 $2|\mathcal{D}|$ 个约束。求出的最优速度曲线满足 $(s^{*}(t_j), t_j) \in \Omega$ 对所有 $j=0,1,\ldots,K$ 成立。工程上走廊约束以 \texttt{STDrivableBoundary} 这一 ST 可行驶边界结构承载，由 \texttt{LaneFollowPath} 的 \texttt{BuildMikuSTDrivableBoundary} 依据路径阶段的时变投影构建，挂载于 \texttt{ReferenceLineInfo} 上下文随数据总线传至 \texttt{PiecewiseJerkSpeedOptimizer}，后者据此按式\eqref{eq:tighten_ub}收紧上界。

走廊上界与 \texttt{SpeedBoundsDecider} 的原有输出在 QP 层取最小值完成双源融合

\begin{equation}
    s_j^{ub} = \min\!\left(s_j^{ub,\mathrm{SBD}},\ s_j^{ub,\mathrm{corridor}}\right), \quad j = 0,1,\ldots,K.
    \label{eq:dual_source}
\end{equation}

两个约束源信息互补。\texttt{SpeedBoundsDecider} 基于已确定路径的碰撞检测描述各障碍物在 path 上的到达时刻；走廊上界来自路径阶段对动态障碍物预测轨迹的分析，描述 path 决策所依赖的时序前提。仅含静态障碍物时 $\mathcal{T}=\varnothing$，走廊约束不生效，由 \texttt{SpeedBoundsDecider} 单独起作用；含动态障碍物时，走廊约束给出精确的等待截止时间 $\tau(s_k)$，把 \texttt{SpeedDecider} 过度保守的全时段让行替换为精确时间窗；动静混合场景中两源各自主导不同时间步，取最小值后 QP 自动得到全局一致的速度曲线。这一融合正好补上路径决策信息未传递至速度规划的设计缺口。

整条传递通道由两条互补路径实现。隐式通道在 SL 阶段产生，当某 $s_k$ 处 $g_{p^{*}}(s_k)<W_{ego}$ 时，按 Apollo 的 blocked-then-trim 机制 path 直接截断在 $s_k$，下游 \texttt{SpeedBoundsDecider} 在被截短的 path 上构建 ST 图，超出 $s_k$ 的纵向区间根本不进入 QP 决策变量，禁行切片信息以 path 几何长度的方式隐式传达。显式通道把 $\mathcal{T}=\{(s_k, \tau_k)\}$ 作为 \texttt{PathBoundsDecider} 的新增输出，经上述 \texttt{STDrivableBoundary} 传递并按式\eqref{eq:tighten_ub}收紧上界。改动范围限于两处，\texttt{PathBoundsDecider} 新增时变 SL 投影与 $\mathcal{T}$ 输出，\texttt{PiecewiseJerkSpeedOptimizer} 新增 $\mathcal{T}$ 读入与上界收紧，而 PathOptimizer 与 \texttt{SpeedBoundsDecider} 不作修改，符合 Apollo 的插件式扩展规范。

$\tau(s)$ 用于构造时变通行带，通行带约束影响速度规划，速度规划结果又决定自车实际到达时间，三者构成循环依赖。本文借 Apollo 每个规划周期重规划的机制以单步延迟隐式处理，每个周期用上周期速度输出更新 $\tau(s)$，在连续多个周期中逐步逼近真实到达时间，该做法与解耦框架现有惯例一致。循环依赖的理论收敛性本文未予证明，作为开放问题列出。

\subsection{整体流程与复杂度}
\label{sec:method_algorithm}

三个组件在一次路径边界求解中顺序衔接。算法\ref{alg:miku}给出从自车状态与障碍物预测轨迹到路径边界与速度约束清单的完整流程。第一步以式\eqref{eq:arrival_time}估计到达时间；第二步据 $\tau(s)$ 做时变 SL 投影并按式\eqref{eq:dynamic_delta}、式\eqref{eq:dynamic_uv}合成差异化裕度与安全边界；第三步扫描线分组；第四步按引理\ref{lem:boundary_trim}做越界预处理；第五步组内按 $u_i$ 升序排列；第六步按式\eqref{eq:gap_def}计算 $k+1$ 个间隙；第七步选最大间隙分配避让方向；第八步输出路径边界与时间清单 $\mathcal{T}$，后者经式\eqref{eq:tighten_ub}注入速度 QP。

\begin{algorithm}[htbp]
\caption{MIKU 多障碍物统一规划求解}
\label{alg:miku}
\KwIn{$n$ 个障碍物及其预测轨迹，道路边界 $[l_{road}^{-}(s), l_{road}^{+}(s)]$，自车状态 $(s_{ego}, v_{ego}, a_{ego})$}
\KwOut{路径边界 $\{[l^{-}(s), l^{+}(s)]\}$，速度约束清单 $\mathcal{T}$}
\For{每个路径离散位置 $s$}{
  按式\eqref{eq:arrival_time}估计到达时间 $\tau(s)$\;
  \For{每个障碍物 $i$}{
    在 $\tau(s)$ 时刻查询预测轨迹并投影至 Frenet，得时变 SL 边界；越出活跃范围则跳过\;
    按式\eqref{eq:threat_score}算威胁度 $\Theta_i$，按式\eqref{eq:dynamic_delta}合成 $\delta_i$，得 $u_i(s)$、$v_i(s)$\;
  }
  按 $s_i^{-}$ 升序扫描线分组，将活跃障碍物划入连通分量\;
  \For{每个连通分量 $C$}{
    按引理\ref{lem:boundary_trim}对越界障碍物做规则预处理\;
    组内按 $u_i$ 升序排列为 $O_{(1)}, \ldots, O_{(k)}$\;
    按式\eqref{eq:gap_def}计算 $k+1$ 个间隙 $g_0, \ldots, g_k$\;
    取 $p^{*} = \arg\max_p g_p$，分配 $d_{(i)}=L\ (i\leq p^{*})$ 或 $R\ (i>p^{*})$\;
  }
  由分割结果写出该截面的 $[l^{-}(s), l^{+}(s)]$，记录禁行切片\;
}
对每个动态障碍物取关键位置 $(s_k, \tau_k)$ 汇入 $\mathcal{T}$\;
\KwRet 路径边界与 $\mathcal{T}$
\end{algorithm}

复杂度可逐段相加。到达时间估计 $O(K)$，$K$ 为路径离散点数；时变投影 $O(n N_s)$，$N_s$ 为单障碍物时变投影采样数；扫描线排序 $O(n\log n)$、扫描 $O(n)$；组内排序合计 $\sum_j O(k_j\log k_j)\leq O(n\log n)$，由 $k\log k$ 的凸性与 $\sum_j k_j = n$ 保证；间隙计算与边界构建 $O(K)$。整体为 $O(n\log n + n N_s + K)$，当 $N_s=O(1)$ 时退化为 $O(n\log n + K)$，$K$ 为路径或时间离散规模。相比之下 Baseline 逐障碍物贪心为 $O(nK)$。单截面层面，MIKU 将暴力枚举的 $O(2^k)$ 降为 $O(k\log k)$，在满足 Apollo 实时预算的同时把方向组合的搜索做到完备。三个组件由此合成一条计算代价几乎不增的显式约束传递通道，第\ref{sec:experiment}节将在仿真场景上验证其通行能力与开销。

% ====== _parts/theory.tex ======
% !TeX root = ../../paper.tex
% 理论保证章节。定理环境若已在导言区定义则不重复定义。
\makeatletter
\@ifundefined{c@lemma}{\theoremstyle{definition}\newtheorem{lemma}{引理}}{}
\@ifundefined{c@theorem}{\theoremstyle{definition}\newtheorem{theorem}{定理}}{}
\@ifundefined{c@proposition}{\theoremstyle{definition}\newtheorem{proposition}{命题}}{}
\makeatother

\section{理论保证}
\label{sec:theory}

本节为前文构建的约束传递通道给出可证明的理论支撑，收敛到两条相互衔接的结论。其一是纯解耦架构的不可行性充分条件，刻画路径阶段在丢失时间维度后报告阻塞、而引入到达时间后联合问题仍可行的情形，由此说明显式传递时间约束的必要性。其二围绕组级最优通行带展开，先以连续性引理把单截面的方向组合搜索空间由 $2^k$ 压缩到 $k+1$ 个候选分割，再以最大间隙最优性定理给出通行带宽度与候选间隙的精确等式，使最优解可经一次线性扫描定位，单截面复杂度由 $O(2^k)$ 降至 $O(k\log k)$，越界障碍物的预处理规则作为定理前提一并给出。最后给出整体复杂度 $O(n\log n + K)$，并证明走廊非空时增强后的速度二次规划仍保有可行解。全节的最优性陈述均限定在所声明的前提之内，凸性来源单独说明，不作超出可证范围的声称。

\subsection{纯解耦的不可行性充分条件}
\label{sec:thy_infeasible}

路径与速度解耦把联合规划拆分为两个顺序子问题。路径子问题在 Frenet 坐标的 $SL$ 平面上确定横向边界 $[l^{-}(s), l^{+}(s)]$，速度子问题在固定路径上构建 $ST$ 图并求解速度曲线。两层之间的接口是确定性的路径几何，时间维度在此被截断。在生产实现中，PathBoundsDecider 对静态障碍物仅采用 $t{=}0$ 时刻的空间快照构建边界，对动态障碍物则经 IsWithinPathDeciderScopeObstacle 中的 IsStatic 判定排除于路径决策之外。无论哪种处理，路径阶段都无从查询障碍物随时间移动后的真实占用，其可用信息至多等价于对预测时域内所有时刻取并集投影，即取障碍物占据的最保守包络。

记路径可行域为 $\mathcal{F}_{\text{dec}}$，时间感知的联合可行域为 $\mathcal{F}_{\text{joint}}$。由于解耦的路径层不修改时间、速度层不修改横向路径，解耦可行域是路径可行域与速度可行域的笛卡尔积；联合可行域则允许时变约束，是更大的集合，故 $\mathcal{F}_{\text{dec}} \subseteq \mathcal{F}_{\text{joint}}$。下述命题刻画该包含关系在何种条件下成为严格包含。

\begin{proposition}[纯解耦不可行性的充分条件]
\label{prop:decoupled_infeasible}
设动态障碍物 $o_i$ 在预测时域 $[0, T_{\text{pred}}]$ 内的 $SL$ 投影为 $\mathrm{Proj}_{SL}(o_i, t)$，自车宽度为 $W_{\text{ego}}$，道路横向边界为 $[l_{road}^{-}(s), l_{road}^{+}(s)]$，记去除半车宽后的可行区间为 $\mathcal{B}_{W}(s) = \bigl[\, l_{road}^{-}(s) + W_{\text{ego}}/2,\ l_{road}^{+}(s) - W_{\text{ego}}/2 \,\bigr]$。若存在纵向区间 $[s_a, s_b]$，使得
\begin{equation}
    \bigcup_{t \in [0, T_{\text{pred}}]} \mathrm{Proj}_{SL}(o_i, t)\ \supseteq\ \mathcal{B}_{W}(s), \qquad \forall\, s \in [s_a, s_b],
    \label{eq:thy_infeasible_cond}
\end{equation}
则时间盲的路径子问题在 $[s_a, s_b]$ 上不存在宽度为 $W_{\text{ego}}$ 的可行通行带，即 $\mathcal{F}_{\text{dec}} = \emptyset$，路径阶段报告 BLOCKED。该条件是解耦失效的充分条件，而非联合问题不可行的必要条件。
\end{proposition}

证明思路。时间盲的路径子问题只能依据并集投影构造占用区域，其在位置 $s$ 处观测到的可用通行带宽度为道路宽度减去并集投影在道路内的横向测度。当式 \eqref{eq:thy_infeasible_cond} 成立时，并集投影覆盖了去除半车宽后的整段可行区间，残余宽度对所有 $s \in [s_a, s_b]$ 均小于 $W_{\text{ego}}$，于是 $l^{+}(s) < l^{-}(s)$，PiecewiseJerkPathOptimizer 无法在其中放置宽度为 $W_{\text{ego}}$ 的横向偏移曲线，路径阶段判定 BLOCKED，触发 TrimPathBounds 截断路径并回退至 FallbackPath 停车。反之，若路径层可利用到达时间函数 $\tau(s)$，则只需查询单一时刻的占用 $\mathrm{Proj}_{SL}(o_i, \tau(s)) \subseteq \bigcup_t \mathrm{Proj}_{SL}(o_i, t)$。当障碍物在自车实际到达时刻 $\tau(s)$ 已移离该横向区间时，残余宽度恢复至不小于 $W_{\text{ego}}$，联合问题在 $[s_a, s_b]$ 上可行。因此并集投影下的不可行并不蕴含联合问题不可行，二者的差正是被解耦接口截断的时间信息所致。$\hfill\Box$

需要强调，$\tau(s)$ 由上周期速度曲线查表或二阶运动学模型估计给定，在路径层是已知输入而非优化变量，命题 \ref{prop:decoupled_infeasible} 因此不涉及对 $\tau(s)$ 的联合优化。该命题揭示解耦失效的根源在于时间维度缺失导致占用区域被过度膨胀，由此引出在解耦接口处显式传递时间约束的核心思路。它给出的是充分条件：满足式 \eqref{eq:thy_infeasible_cond} 必然触发解耦阻塞，但解耦阻塞也可能源于后文讨论的逐障碍物贪心决策的虚假阻塞，二者需加区分。

\subsection{组级最优通行带的最优性}
\label{sec:thy_max_gap}

考虑单个 $s$ 截面上同时活跃的 $k$ 个障碍物。每个障碍物 $o_i$ 经差异化安全裕度 $\delta_i$ 膨胀后得到右侧安全边界 $u_i = l_i^{-} - \delta_i$ 与左侧安全边界 $v_i = l_i^{+} + \delta_i$，避让决策变量 $d_i \in \{L, R\}$ 二值取值。给定方向向量 $\mathbf{d}$ 后，路径上界取所有右绕障碍物右边界与道路左边界的最小值，下界取所有左绕障碍物左边界与道路右边界的最大值，通行带宽度 $W(\mathbf{d}) = l^{+}(\mathbf{d}) - l^{-}(\mathbf{d})$。逐障碍物独立赋值的朴素搜索空间为 $2^k$。逐障碍物贪心仅检查其中一种组合，在多数场景下能给出可行解，但在特定对称或密集构型下可能选中分裂式组合而遗漏可行解。两个分居参考线两侧且近似对称的障碍物即为典型情形，四种方向组合中两种分裂式赋值会从两侧同时收缩边界、把自车挤入中间窄缝而判定 BLOCKED，另两种统一绕行赋值则可借道路另一侧的充裕空间通过。下述引理与定理给出全局最优解的精确刻画。

\begin{lemma}[最优分割的连续性]
\label{lem:continuous_partition}
将障碍物按 $u_i$ 值升序排列后，存在最优解 $\mathbf{d}^{*}$，使得分割在排序中是连续的，即存在 $p \in \{0, 1, \ldots, k\}$ 满足
\begin{equation}
    d_{(i)}^{*} = L,\ \forall\, i \leq p; \qquad d_{(i)}^{*} = R,\ \forall\, i > p.
    \label{eq:continuous_partition}
\end{equation}
\end{lemma}

证明思路。任取一最优解 $\mathbf{d}^{*}$，记其左绕集合为 $\mathcal{L}$、右绕集合为 $\mathcal{R}$。若 $\mathcal{R}=\emptyset$ 或 $\mathcal{L}=\emptyset$，分别对应 $p=k$ 或 $p=0$，结论显然成立。两侧均非空时，令 $\beta = \min_{i \in \mathcal{R}} u_i$，将所有满足 $u_i \geq \beta$ 且原属 $\mathcal{L}$ 的障碍物改入 $\mathcal{R}$，其余决策不变。新加入 $\mathcal{R}$ 的障碍物均有 $u_i \geq \beta$，故 $\min_{i \in \mathcal{R}} u_i$ 不变，上界 $l^{+}$ 不减；$\mathcal{L}$ 只减少元素，故 $\max_{i \in \mathcal{L}} v_i$ 不增，下界 $l^{-}$ 不增，若 $\mathcal{L}$ 被移空则下界退化为 $l_{road}^{-}$。因此调整后通行带宽度不小于原最优解，仍为最优解。此时所有 $u_i < \beta$ 的障碍物在 $\mathcal{L}$、所有 $u_i \geq \beta$ 的障碍物在 $\mathcal{R}$，按 $u_i$ 升序必存在分割点 $p$ 使前 $p$ 个属于 $\mathcal{L}$、后 $k-p$ 个属于 $\mathcal{R}$，即式 \eqref{eq:continuous_partition} 成立。$\hfill\Box$

引理 \ref{lem:continuous_partition} 把搜索空间由 $2^k$ 个任意分配降为 $k+1$ 个候选分割点 $p = 0, 1, \ldots, k$。其中 $p=0$ 表示全部右绕即自车走最右侧通道，$p=k$ 表示全部左绕即自车走最左侧通道，$1 \leq p \leq k-1$ 表示自车从第 $p$ 与第 $p+1$ 个障碍物之间的间隙通过。引理仅依赖 $u$ 值的排序性，不要求 $\delta_i$ 为常数；当存在 $u_i$ 相等的障碍物时，可取 $v_i$ 降序为二级排序键以保证排序唯一，不影响结论。

定理的精确等式以两条道路边界约束为前提。当某障碍物膨胀后越出道路边界时，需先经预处理使前提成立。

\begin{lemma}[越界障碍物的预处理规则]
\label{lem:boundary_trim}
若存在障碍物 $i$ 使得 $u_i < l_{road}^{-}$ 或 $v_i > l_{road}^{+}$，则按以下规则预处理后，定理 \ref{thm:max_gap} 的结论仍然成立。
\begin{enumerate}
    \item 规则 A：若 $u_i < l_{road}^{-}$ 且 $v_i > l_{road}^{+}$，障碍物膨胀后完全覆盖道路横截面，判定当前 $s$ 截面 BLOCKED，不进入后续分割搜索。
    \item 规则 B：若 $u_i < l_{road}^{-}$ 且 $v_i \leq l_{road}^{+}$，障碍物右膨胀区越过道路右边界，强制 $d_i = L$ 并从候选列表移除，其约束由 $v_i$ 并入 $\max$ 运算，不影响其它分割点的间隙。
    \item 规则 C：若 $u_i \geq l_{road}^{-}$ 且 $v_i > l_{road}^{+}$，障碍物左膨胀区越过道路左边界，强制 $d_i = R$ 并从候选列表移除，其约束由 $u_i$ 并入 $\min$ 运算，不影响其它分割点的间隙。
\end{enumerate}
\end{lemma}

证明思路。规则 A 下 $l^{+}(\mathbf{d}) < l^{-}(\mathbf{d})$ 对任意 $\mathbf{d}$ 成立，无可行分割，等价于最优宽度 $W^{*}=0$。规则 B 中若分配 $d_i = R$，则 $l^{+}(\mathbf{d}) \leq u_i < l_{road}^{-} \leq l^{-}(\mathbf{d})$，通行带宽度非正，故 $d_i = L$ 为唯一最优选择；强制左绕后 $v_i$ 仅参与左安全边界的聚合，不改变其余障碍物之间的间隙序。规则 C 与规则 B 对称。三条规则预处理后，剩余障碍物全部满足 $u_j \geq l_{road}^{-}$ 且 $v_j \leq l_{road}^{+}$，定理 \ref{thm:max_gap} 的前提成立。$\hfill\Box$

基于引理 \ref{lem:continuous_partition}，只需枚举 $k+1$ 个候选分割点。对分割点 $p$，由于仅按 $u_i$ 排序不能保证 $v_i$ 单调，需先定义前缀左安全边界
\begin{equation}
    V_p = \max_{1 \leq q \leq p} v_{(q)}, \qquad p = 1, \ldots, k,
    \label{eq:prefix_v}
\end{equation}
并定义包含两个端部在内的 $k+1$ 个间隙
\begin{equation}
    g_p = \begin{cases}
        u_{(1)} - l_{road}^{-} & p = 0, \\[0.3em]
        u_{(p+1)} - V_p & p = 1, 2, \ldots, k-1, \\[0.3em]
        l_{road}^{+} - V_k & p = k.
    \end{cases}
    \label{eq:gap_def}
\end{equation}
其中 $g_0$ 为道路右边界到最右侧障碍物之间的间隙，$g_k$ 为前 $k$ 个障碍物的左安全包络到道路左边界之间的间隙，中间各 $g_p$ 为前 $p$ 个障碍物的最左安全边界 $V_p$ 与后续障碍物最右安全边界 $u_{(p+1)}$ 之间的间隙。当障碍物横向边界不交错时 $V_p = v_{(p)}$，该定义退化为相邻障碍物之间的普通几何间隙。

\begin{theorem}[最大间隙最优性]
\label{thm:max_gap}
若所有障碍物的安全裕度扩展后均未超出道路边界，即对任意 $i$ 均有 $u_i \geq l_{road}^{-}$ 且 $v_i \leq l_{road}^{+}$，则对分割点 $p$，通行带宽度 $W(p) = g_p$，最优通行带宽度为
\begin{equation}
    W^{*} = \max_{p = 0, 1, \ldots, k} g_p.
    \label{eq:max_gap_opt}
\end{equation}
\end{theorem}

证明思路。对各分割点直接代入验证。当 $p = 0$ 即全部右绕时，$l^{+}(0) = \min(l_{road}^{+}, u_{(1)}) = u_{(1)}$、$l^{-}(0) = l_{road}^{-}$，故 $W(0) = g_0$。当 $1 \leq p \leq k-1$ 时，由前提 $u_{(p+1)} \leq l_{road}^{+}$ 得 $l^{+}(p) = u_{(p+1)}$，由前提 $V_p \geq l_{road}^{-}$ 得 $l^{-}(p) = V_p$，故 $W(p) = u_{(p+1)} - V_p = g_p$。当 $p = k$ 即全部左绕时，$l^{+}(k) = l_{road}^{+}$、$l^{-}(k) = V_k$，故 $W(k) = g_k$。化简仅用到前提 $u_i \geq l_{road}^{-}$、$v_i \leq l_{road}^{+}$ 与安全扩展区间的基本关系 $u_i \leq v_i$。因此最大 $g_p$ 对应的分割点 $p^{*}$ 即最大化 $W(p)$ 的分割点，式 \eqref{eq:max_gap_opt} 严格成立。$\hfill\Box$

定理 \ref{thm:max_gap} 把通行带宽度与候选间隙建立精确等式，使最优分割可经一次线性扫描定位。三类边界情形需单独处理。其一，当安全包络交错即 $V_p > u_{(p+1)}$ 时，对应间隙 $g_p < 0$，表示通道不可通行，算法自然跳过。其二，当全部 $k+1$ 个间隙均不足车宽，即 $\max_p g_p < W_{\text{ego}}$ 时，该截面无可通行路径，这是真实的 BLOCKED 而非算法局限，对应命题 \ref{prop:decoupled_infeasible} 中并集投影外仍不可行的情形。其三，引理 \ref{lem:continuous_partition} 的连续性使最优解必落在 $k+1$ 个候选分割之中，与逐障碍物贪心仅检查单一组合相比，最大间隙搜索遍历全部候选，因而在对称或密集构型下不会输出虚假 BLOCKED；唯有当全部间隙确实为负时输出的 BLOCKED 才是真实阻塞。

\subsection{整体复杂度}
\label{sec:thy_complexity}

单截面层面，引理 \ref{lem:continuous_partition} 把搜索空间由 $2^k$ 压缩到 $k+1$ 个候选分割，定理 \ref{thm:max_gap} 进一步将每个候选的宽度评估归约为间隙查表。组内按 $u_i$ 排序需 $O(k \log k)$，越界预处理与前缀最大值递推各为 $O(k)$，最大间隙扫描为 $O(k)$，故单截面复杂度由暴力枚举的 $O(2^k)$ 降至 $O(k \log k)$。

将单截面方法推广至全路径需先把 $n$ 个障碍物按纵向重叠关系分组。两障碍物纵向重叠当且仅当其纵向区间相交，该关系对称但不传递。取其传递闭包将障碍物划分为若干连通分量，不同分量的纵向区间互不相交，任一 $s$ 截面至多受一个分量影响，分量间的避让决策可独立求解。扫描线分组按 $s_i^{-}$ 升序排列后单遍扫描合并，排序为 $O(n \log n)$、扫描为 $O(n)$。各连通分量的组内排序之和满足 $\sum_j O(k_j \log k_j) \leq O(n \log n)$，该上界由 $k \log k$ 的凸性与 $\sum_j k_j = n$ 保证。沿路径离散点构建边界并输出速度约束清单为 $O(K)$，$K$ 为路径与时间离散规模。综合各阶段，整体复杂度为 $O(n \log n + K)$，相对 Baseline 逐障碍物处理的 $O(nK)$ 在障碍物数量增长时具有更优的可扩展性，且在 Apollo 11.0 的 $100$~ms 规划周期约束下满足实时性要求。

链式连通分量需补充说明。连通分量不保证存在所有障碍物共同活跃的公共截面，链式情形中 $o_1$ 与 $o_2$ 重叠、$o_2$ 与 $o_3$ 重叠，而 $o_1$ 与 $o_3$ 可完全不重叠。对此采用保守包络处理，把同一分量内所有障碍物视为同时施加约束。由于任一截面实际活跃的障碍物集合是分量的子集，包络宽度是实际可用宽度的下界，包络宽度不小于车宽是该分量可行的充分条件。包络不可行仅说明充分条件失败，不意味着真实场景必然不可行；保守包络保证安全性，但不对一般链式分量声称给出全局最优解。

\subsection{走廊非空时的可行性恢复}
\label{sec:thy_feasibility}

前文将 $SL$ 平面的逐截面通行性沿到达时间 $\tau(s)$ 抬升至 $ST$ 平面，得到时空可通行走廊 $\Omega$，并经 STDrivableBoundary 以收紧上界的方式注入速度二次规划。本小节说明该注入在何种前提下保持速度问题的可行性。

走廊 $\Omega$ 由 $\tau(s)$ 的左下沿与若干水平禁行切片组合而成，其中禁行切片对应某纵向位置处最优间隙不足车宽时整列被排除的情形，因此 $\Omega$ 作为集合一般是非凸的。凸性来源不在 $\Omega$ 本身，而在其沿离散时间步的逐时间步线性投影。注入过程对每个关键位置仅按到达时间对相应时间步的纵向上界做一次取最小值的收紧，即把走廊约束化为速度变量的盒式上界，与 SpeedBoundsDecider 给出的 $ST$ 边界取最小值融合。该盒式约束与 Baseline 在 PathTimeHeuristicOptimizer 给出让行或超车组合后所得的线性约束在二次规划结构上一致，因此 PiecewiseJerkSpeedOptimizer 的目标矩阵、等式约束与稀疏模式完全不变，仅上界向量的若干分量被收紧。

\begin{proposition}[走廊非空时的可行性恢复]
\label{prop:feasibility_recovery}
设速度二次规划的目标为连续二次凸函数，原始约束为 SpeedBoundsDecider 给出的 $ST$ 盒式边界。若走廊约束注入后，在速度二次规划的离散时间网格上其可行域非空且仍为凸集，则增强后的速度二次规划必有最优解。
\end{proposition}

证明思路。走廊约束以逐时间步取最小值的方式收紧纵向上界，每个时间步上得到的仍是闭区间盒式约束，多个盒式约束的交集为凸集。当该交集与原始 $ST$ 盒式边界的公共部分非空时，增强后的可行域是有界闭凸集。连续二次凸函数在非空有界闭凸集上必取得最小值，故增强后的速度二次规划存在最优解。$\hfill\Box$

命题 \ref{prop:feasibility_recovery} 的前提需明确界定。其一，结论限定在离散时间网格上，即逐时间步线性投影后的有限维盒式约束，而非对非凸集合 $\Omega$ 整体声称凸性。其二，可行域非空是前提而非结论，当障碍物在自车到达前无法让出足够空间时收紧后的上界可能与下界冲突，此时走廊为空，对应真实需要等待或停车的物理情形，不属于可恢复范畴。其三，走廊由 $\tau(s)$ 构造，而 $\tau(s)$ 又依赖速度规划结果，二者构成循环依赖。本文借助 Apollo 每个规划周期的周期性重规划以单步延迟方式处理，每周期用上周期速度输出更新 $\tau(s)$，与解耦框架的现有做法一致；该近似的理论收敛性分析作为开放问题，本文不予声称。

把命题 \ref{prop:decoupled_infeasible} 与命题 \ref{prop:feasibility_recovery} 合并来看，二者刻画了约束传递通道两端的理论效果。前者表明纯解耦在并集投影下可能报告不可行，而联合问题在引入到达时间后仍可行，说明显式传递时间约束的必要性；后者表明在走廊非空且离散网格上凸的前提下，注入约束不破坏速度问题的可解性，说明该传递在计算上是良性的。中间环节由引理 \ref{lem:continuous_partition}、引理 \ref{lem:boundary_trim} 与定理 \ref{thm:max_gap} 保证空间维度的全局最优分割可在 $O(k \log k)$ 内求得，整体复杂度 $O(n \log n + K)$ 满足实时约束。上述结论均限定在各自声明的前提之内，凸性来源与循环依赖等局限已逐条标注。

% ====== _parts/experiment.tex ======
\section{实验与分析}
\label{sec:experiment}

本节在Apollo 11.0生产级代码的工程接入之上，用数值仿真量化显式约束传递通道带来的通行能力、轨迹质量与求解开销变化。实验分为四部分，依次给出八场景的通行成功率与速度对照、可比场景下的求解效率对照、五个机制的组件消融，以及威胁度权重的灵敏度分析。本节末尾单列局限性，主动声明本工作在验证边界上的不足。

连续数值指标取自一套以Python复现Apollo Planning主要计算环节的仿真程序，依次模拟路径与速度规划各阶段，二次规划子问题统一交由OSQP求解；Baseline与MIKU在同一程序内并列实现，对每个场景同步输出SL路径边界、ST走廊与纵向速度、纵向加速度、横向加速度、jerk四条时序曲线。Dreamview在环截图来自MIKU的C++实现在Apollo工程中的运行，仅用于展示接入后的规划行为，不与上述数值指标混作同一数据源。

\subsection{仿真平台与场景设置}
\label{subsec:exp_setup}

实验在桌面级算法验证平台上进行，处理器为Intel i9-14900HX，内存64\,GB DDR5，操作系统为ArchLinux；软件栈为Apollo 11.0、C++17、OSQP 0.6.2与Dreamview 2.0，数值复现程序基于Python 3.12。需要强调，本平台为算法验证用途，QP求解耗时等指标仅反映该桌面环境的表现，车载嵌入式平台的计算性能须另行评估，本工作不就此外推。

MIKU的关键参数沿用威胁度建模与通行带构造两节的标定取值。差异化安全裕度$\delta_i$取值区间为$[0.10, 0.40]$\,m，障碍物聚类半径为10\,m，纵向重叠分组阈值为12\,m。这些取值面向城市道路工况，未针对其他工况重新标定。

测试场景共八个，由四个原始压力场景与四个配对可比场景构成。压力场景P1--P4保留临界构型，用于暴露Baseline在动态障碍物、统一安全裕度与逐障碍物贪心决策三类配置下的失效边界。可比场景C1--C4逐一复制对应压力场景，仅轻量调整规划时间窗或障碍物横向位置，使Baseline与MIKU均能完成通过，从而在同一可达前提下比较轨迹质量与求解开销。如此设计既保留了原构型对各机制必要性的验证，也避免把通过与失败的离散对照直接解读为连续指标的优势。八个场景的构型与可比化调整如表\ref{tab:exp_scenarios}所示。

\begin{table}[htbp]
\centering
\caption{压力场景与配对可比场景的构型设置}
\label{tab:exp_scenarios}
\small
\begin{tabular}{c l p{7.4cm}}
\toprule
\textbf{场景对} & \textbf{构型} & \textbf{可比化调整与比较目的} \\
\midrule
P1 / C1 & 单行人横穿 & C1延长规划时间窗，比较动态障碍物提前投影后的纵向平顺性 \\
P2 / C2 & 行人加停车 & C2外移停车横向中心使Baseline恢复可达，比较动静混合分组后的速度QP开销 \\
P3 / C3 & 窄路加双侧交通锥 & C3将对称锥列改为轻微非对称使Baseline可通过，保留差异化裕度释放通行空间 \\
P4 / C4 & 单车道维修封闭加借道 & C4将水马横向中心调至Baseline贪心切换临界侧，比较组级边界组织的求解开销 \\
\bottomrule
\end{tabular}
\end{table}

这里需要界定Baseline失败的比较口径。Apollo原生的PathBoundsDecider通过IsWithinPathDeciderScopeObstacle将动态障碍物排除在路径阶段之外，交由SpeedBoundsDecider在ST阶段以时间窗约束兜底；MIKU则通过时变SL投影把动态障碍物显式纳入路径边界构建。因此P1--P4中Baseline的失败并不等于Apollo整套流水线的失败，而是特指路径阶段单独处理动态障碍物这一能力口径下的失效复现。可比场景在同一口径下放宽时间窗或障碍物位置，使Baseline的ST阶段兜底可通行，以便在可达前提下对照两种方案的连续指标。

每个场景计算七项指标，覆盖通行能力、轨迹舒适性与求解开销三类。通行能力由通过率$P_{\mathrm{pass}}$与平均速度$\bar v$度量，其中$P_{\mathrm{pass}}$以仿真终点满足$s_{\mathrm{end}} \geq s_{\max} - 1$\,m为通过判据。轨迹舒适性由最大纵向加速度幅值$|a_x|_{\max}$、最大横向加速度幅值$|a_y|_{\max}$与最大jerk幅值$|j|_{\max}$度量，其中横向加速度由路径曲率诱导，$a_y = v^2 \kappa$。求解开销由路径QP与速度QP合计的总耗时$t_{\mathrm{qp}}$度量。后续各结果表均沿用这一口径。

\subsection{通行成功率与速度对照}
\label{subsec:exp_passrate}

八场景的通行成功率与平均速度对照如表\ref{tab:exp_pass}所示。在压力场景P1--P4中，Baseline通过0/4，MIKU通过4/4；显式约束传递通道使路径阶段获得时间感知能力，避免了Baseline把动态障碍物按$t=0$快照过度膨胀后强制制动的失效。在可比场景C1--C4中，两者均通过4/4，说明在Baseline本可通行的常规构型下，MIKU不损失通行能力。汇总八场景，Baseline通过4/8，MIKU通过8/8。

\begin{table}[htbp]
\centering
\caption{八场景通行成功率与平均速度对照}
\label{tab:exp_pass}
\small
\begin{tabular}{lccc}
\toprule
\textbf{场景组} & \textbf{通过率 Baseline / MIKU} & \textbf{$\bar v$ Baseline (m/s)} & \textbf{$\bar v$ MIKU (m/s)} \\
\midrule
压力场景 P1--P4 & 0/4 \ / \ 4/4 & 2.58 & 4.66 \\
可比场景 C1--C4 & 4/4 \ / \ 4/4 & 4.48 & 4.47 \\
\midrule
汇总 八场景 & 4/8 \ / \ 8/8 & 3.53 & 4.57 \\
\bottomrule
\end{tabular}
\end{table}

汇总平均速度由Baseline的3.53\,m/s提升至MIKU的4.57\,m/s，相对提升29.3\%。这一提升须谨慎解读。可比场景中两者速度分别为4.48\,m/s与4.47\,m/s，基本持平，说明在双方均可通行的构型下速度并非MIKU的优势所在。汇总速度的差距主要来自压力场景，Baseline在该组触发停车决策使平均速度被拉低至2.58\,m/s，而MIKU维持4.66\,m/s。因此29.3\%反映的是压力场景中Baseline失效拖低均值后的汇总效应，而非常规巡航速度的普遍提升，本工作不将其解释为巡航性能的整体改善。

可比场景的轨迹质量对照进一步说明约束前置的作用。MIKU把避让决策提前至路径阶段后，速度阶段不再需要被动急刹，C1--C2动态场景下的纵向加速度峰值与jerk峰值相应低于Baseline，纵向行驶更平顺；C3窄路场景下两者横向加速度峰值均接近零，差异化裕度并未以牺牲横向舒适性为代价换取通行空间。

\subsection{求解效率对照}
\label{subsec:exp_qp}

QP求解总耗时的对照如表\ref{tab:exp_qp}所示。需要先说明效率比较的口径。在压力场景P1--P4中Baseline进入FallbackPath降级分支，QP求解提前中止于减速停车阶段，记录的耗时反映降级前的不完整计算，不构成两种算法的效率比较依据。因此本工作的求解效率结论仅基于C1--C4可比场景，该组两者均完成完整求解，耗时可比。

\begin{table}[htbp]
\centering
\caption{QP求解平均总耗时对照（ms）}
\label{tab:exp_qp}
\small
\begin{tabular}{lcc}
\toprule
\textbf{场景组} & \textbf{Baseline} & \textbf{MIKU} \\
\midrule
压力场景 P1--P4（仅参考，Baseline降级） & 23.96 & 4.86 \\
可比场景 C1--C4（效率结论依据） & 19.16 & 6.26 \\
\midrule
汇总 八场景 & 21.56 & 5.56 \\
\bottomrule
\end{tabular}
\end{table}

在可比场景C1--C4中，Baseline与MIKU的平均QP总耗时分别为19.16\,ms与6.26\,ms。MIKU将部分决策判断前移至路径阶段后，速度QP的约束规模下降，路径阶段新增的扫描线分组与时变投影开销保持在毫秒量级，可比场景的总体求解效率随之改善，且通行能力未受影响。从复杂度看，MIKU新增的分组与投影步骤为$O(n \log n)$，单截面方向决策由穷举的$O(2^k)$降为排序后的$O(k \log k)$，整体复杂度为$O(n \log n + K)$，其中$K$为路径与时间离散规模。走廊约束以标准ST上界形式注入，OSQP对约束数量的增量不敏感，不改变QP的稀疏结构。

\subsection{组件消融}
\label{subsec:exp_ablation}

消融实验只使用P1--P4压力场景，用于定位显式约束传递通道中各机制对通行成功率的边际贡献。本工作把通道分解为五个正交机制，记为C1至C5，分别为时变SL投影、扫描线纵向分组、组内最大间隙决策、多因子差异化裕度$\delta_i$与ST走廊注入。其中扫描线纵向分组与组内最大间隙决策在工程实现中为一体的两步工序，分组输出直接作为间隙搜索的输入，单独关闭其一不构成独立配置，故合并为一个消融变体。采用减法消融，在完整MIKU基础上逐一关闭单个机制，得到六个对照变体，如表\ref{tab:exp_ablation_matrix}所示。

\begin{table}[htbp]
\centering
\caption{消融变体与机制开关}
\label{tab:exp_ablation_matrix}
\small
\begin{tabular}{llccccc}
\toprule
\textbf{编号} & \textbf{变体} & 投影 & 分组 & 间隙 & 裕度 & 走廊 \\
\midrule
M0 & Baseline             & $\times$    & $\times$    & $\times$    & $\times$    & $\times$    \\
M1 & 关闭时变投影         & $\times$    & $\checkmark$ & $\checkmark$ & $\checkmark$ & $\checkmark$ \\
M2 & 关闭分组与间隙决策   & $\checkmark$ & $\times$    & $\times$    & $\checkmark$ & $\checkmark$ \\
M3 & 关闭差异化裕度       & $\checkmark$ & $\checkmark$ & $\checkmark$ & $\times$    & $\checkmark$ \\
M4 & 关闭走廊注入         & $\checkmark$ & $\checkmark$ & $\checkmark$ & $\checkmark$ & $\times$    \\
M5 & 完整MIKU             & $\checkmark$ & $\checkmark$ & $\checkmark$ & $\checkmark$ & $\checkmark$ \\
\bottomrule
\end{tabular}
\end{table}

为把多维指标压缩为单一可比标量，本工作采用两层评分。第一层为安全硬门，要求场景成功通过且未触发blocked，任一项失败则该变体在该场景直接计零分，避免短板被高分项稀释。第二层把软指标按四个维度分组，组内子指标等权乘积，组间按权重加权几何平均，得到$0$至$1$区间的场景得分。四个维度的权重为通行效率0.30、轨迹平顺0.30、决策鲁棒0.25、计算开销0.15。各子指标在场景内以完整MIKU作分母做相对归一化，选其而非Baseline作参考，是因为Baseline在多个场景触发blocked导致部分指标不收敛，而完整MIKU为已验证可行解，作分母数值稳定。单场景得分由硬门系数$G \in \{0, 1\}$与四维加权几何平均共同构成
\begin{equation}
S_{\mathrm{scn}} = G \cdot \prod_{d=1}^{4} \left[ \prod_{i \in d} n_i^{1/|d|} \right]^{w_d},
\label{eq:exp_score_scn}
\end{equation}
其中$n_i$为子指标的归一化值，$w_d$为维度权重。全局得分对所有压力场景再做几何平均并放大至百分制，若任一场景因硬门得零分则退回算术平均并折半惩罚。几何平均的否决性使单项失效直接拉低总分，而非被其他高分项掩盖。

六个变体的综合评分如表\ref{tab:exp_ablation_score}所示，以完整MIKU为100分基准。

\begin{table}[htbp]
\centering
\caption{消融变体综合评分（完整MIKU为100基准）}
\label{tab:exp_ablation_score}
\small
\begin{tabular}{llcc}
\toprule
\textbf{编号} & \textbf{变体} & \textbf{通过数} & \textbf{综合评分} \\
\midrule
M5 & 完整MIKU             & 4/4 & 100.00 \\
M4 & 关闭走廊注入         & 4/4 & 99.18 \\
M2 & 关闭分组与间隙决策   & 3/4 & 37.13 \\
M3 & 关闭差异化裕度       & 3/4 & 34.50 \\
M1 & 关闭时变投影         & 2/4 & 25.29 \\
M0 & Baseline             & 0/4 & 0.00 \\
\bottomrule
\end{tabular}
\end{table}

综合评分在很大程度上反映的是通过率差异而非通过场景内的质量差异，这是硬门与几何平均否决性的直接结果。M1、M2、M3因在关键压力场景失败被硬门归零，几何平均随之退化，全局得分均落到38分以下。结合各变体的失效场景可以定位机制的必要性。关闭时变投影的M1在单行人横穿与行人加停车两个动态压力场景失败，动态行人退化为$t=0$快照占据车道，路径阶段被错误约束至当前位置而强制减速等待。这说明一旦选择把动态障碍物纳入路径阶段，就必须配合时变投影，否则反而产生副作用。关闭分组与间隙决策的M2在单车道封闭导流场景失败，该场景的多个静态障碍物构成单一大型连通分量，逐障碍物贪心决策下入口与维持段各障碍物独立选择避让方向，产生左右边界冲突而触发blocked-then-trim停车。关闭差异化裕度的M3在窄路交通锥场景失败，双侧锥间隙仅略大于车身宽度，统一额外缓冲对小尺寸锥桶过度保守，把通行带压缩至不可行。由此，时变投影与分组间隙决策共同决定动态与大型连通分量场景的通行成功率，差异化裕度决定窄路锥列下的可行性。

关闭走廊注入的M4得99.18分，与完整MIKU仅相差0.82分，近似持平。这一结果需要诚实说明，它并非否定走廊注入这一机制。当前四个压力场景均为单车道静态或低动态构型，路径阶段已通过时变投影提前规避了动态时间冲突，速度阶段仅需保守通过，走廊约束未触发活跃边界；关闭它只是少了一组几乎不生效的速度约束。走廊注入的意义在于为路径-速度解耦架构补全此前缺失的约束传递通道，把路径阶段对动态障碍物的时空决策以到达时间形式注入速度QP。在多动态障碍物时间窗交织或交叉路口冲突点决策等场景中，它将成为速度规划的必要约束源。因此走廊注入的定量贡献需要后续在动态冲突更密集的场景中补充验证，本实验数据不足以单独支撑其必要性，其价值由架构设计论证而非本组消融数据给出。

\subsection{灵敏度分析}
\label{subsec:exp_sensitivity}

威胁度评分由五个因子加权构成，权重取值依赖人工标定。为检验通行决策对权重取值的鲁棒性，本工作在P1--P4压力场景上对五个权重施加$\pm 20\%$扰动，做100次蒙特卡洛采样，观察通行带宽度相对基准取值的偏差。结果显示，最大相对偏差为6.4\%，出现在场景四，单因子越界概率为0，即无任何一次扰动使通行带宽度突破工程容忍阈值。

需要单独说明翻转率这一指标。场景四的避让方向排序翻转率最坏达57\%，但这些翻转集中发生在威胁度极为接近的同类交通锥之间，对应的威胁度差异不足0.02，差异化裕度的变化不足0.01\,m，对通行带宽度的影响可以忽略。也就是说，翻转改变的只是两个近乎等价的同类障碍物之间的排序名次，并不改变最终的避让方向决策与通行结果。其余三个场景的排序翻转率均为0。综合来看，权重的合理扰动不改变障碍物之间的相对威胁等级，对通行带宽度的影响处于工程容忍范围内。需要补充的是，上述权重面向城市道路中低速工况标定，在高速大速度差场景下需重新评估。

\subsection{局限性}
\label{subsec:exp_limitations}

本工作主动声明若干验证边界上的不足，以划清结论的适用范围。

本工作未做实车验证。全部连续指标来自Python数值仿真，工程行为展示来自Dreamview在环运行，两者均在仿真层面，定位为生产级代码在环加数值仿真的原型验证，结论尚不能外推到实车。

实验仅在单一桌面级平台上完成。QP求解耗时等开销指标反映Intel i9-14900HX桌面环境的表现，车载嵌入式平台的实时性须另行评估，本工作不就此外推。

压力场景的选取偏向Baseline的失效边界。P1--P4均为Baseline在路径阶段单独处理动态障碍物时失败的临界构型，配对的可比场景C1--C4说明常规可达情形下两者通行能力持平。因此本工作主张的是特定压力构型下的失效复现与恢复，而非对Baseline的普遍碾压。

实验为确定性仿真的单次运行，未做多次重复与统计显著性检验。各场景指标为可复现的单次确定性结果，灵敏度分析中的蒙特卡洛采样仅针对威胁度权重扰动，不构成对全部指标的统计推断。

评价指标未包含最小安全间距等专门的安全性度量。本工作以通过率、速度、加速度、jerk与求解耗时刻画通行与舒适性，安全间距的量化评估列为后续工作。

走廊注入机制在当前场景未被激活。如消融部分所述，四个压力场景的时空决策已在路径阶段由时变投影完成，走廊约束未触发活跃边界，其定量贡献有待在动态冲突更密集的场景中补充验证。

% ====== _parts/related.tex ======
% ============================================================
\section{相关工作}
\label{sec:related}
% ============================================================

路径--速度解耦自~\citet{kant1986pvd}~提出以来一直是结构化道路轨迹规划的主流范式，围绕其展开的改进与替代工作可归纳为三条技术路线。第一条放弃解耦，在时间--纵向--横向构成的三维域内直接求解联合轨迹；第二条保留解耦框架，在子问题内部注入额外信息以缓解信息断裂；第三条即本文路线，在两阶段接口处建立显式约束传递通道。此外，混合整数规划与端到端学习两类方法从组合优化与一体化建模的不同范式给出了思路。下面分别讨论，并逐条说明 MIKU 与各路线的差异化定位。

\subsection{时空联合规划}
\label{sec:related-joint}

时空联合规划放弃路径--速度解耦，转而在三维时空域内直接搜索或优化轨迹，以换取时空一致性。\citet{yang2021_ia_planner}~提出 IA Planner，通过瞬时分析模型仅考察同一时刻的碰撞关系，将三维时空约束投影为二维空间约束后建模为模型预测控制（Model Predictive Control, MPC）问题联合求解，并以启发式速度引导提升纵向运动连续性。\citet{hu2023_hybrid_astar_slt}~在 Frenet 坐标系下构建三维 $s$-$l$-$t$ 时空栅格，先用改进混合 A* 算法粗搜索初始轨迹，再基于初始轨迹构建可行驶时空走廊并以非线性模型预测控制（Nonlinear Model Predictive Control, NMPC）精化；其后续工作~\citep{hu2025_dp_nmpc}~将粗搜索替换为确定性采样的动态规划（Dynamic Programming, DP），以多维时空网格过滤相似节点抑制维度爆炸。\citet{qiao2025stjoint}~则在结构化道路下引入时间--纵向--横向三维联合决策，并与时空安全走廊的二次规划相结合。这类方法的搜索与优化均在三维域内进行，计算开销随障碍物数量上升而显著增加；\citet{hu2023_hybrid_astar_slt}~报告其方法在换道超车与旁车切入两类场景下的平均耗时分别为 117\,ms 与 76\,ms，明显高于传统解耦基线。MIKU 不进入三维联合域，而是在保留两个二维子问题独立求解结构的前提下，仅在接口处补足被解耦丢弃的时间维度信息，从而在不承担联合求解开销的情况下恢复时空协调能力。

\subsection{解耦框架内改进}
\label{sec:related-decoupled}

第二条路线保留解耦框架，通过在子问题内部注入额外信息来缓解信息断裂。\citet{han2026_contingency}~面向自动代客泊车场景提出 Contingency 路径--速度迭代规划器，将动态障碍物的多模态预测概率以软代价形式注入路径二次规划的目标函数，隐式传递预测信息；该方式在目标函数层面引导轨迹偏好，但不对结果施加硬性约束，因而缺乏可行性保证。\citet{zhang2023_dp_rcg}~先以较大时间间隔分别生成候选路径与候选速度曲线集合，再通过静态风险合作博弈在二者间选取最优配对，最后用凸优化精化轨迹；该方法以并行枚举换取路径与速度的协同，组合规模为 $O(mn)$，其中 $m$、$n$ 分别为候选路径数与候选速度数，架构改动量也相应较大。与软代价的隐式传递不同，MIKU 以硬约束形式将空间决策结构化地注入速度子问题，可行性可被刻画并证明；与枚举--博弈不同，MIKU 不构造候选组合的笛卡尔积，而是在单截面将避让方向选择归约为最大间隙问题，单截面复杂度由 $O(2^k)$ 降至 $O(k\log k)$，整体为 $O(n\log n + K)$，且不改变各阶段内部求解结构。本文在理论分析中给出纯解耦规划器不可行性的充分条件，以及通行带与时空走廊非空时的可行性恢复保证。

\subsection{解耦规划的工业应用}
\label{sec:related-industrial}

路径--速度解耦在量产系统中占据主导地位，这也是本文选择在该架构内改进而非将其替换的现实依据。\citet{fan2018baidu}~描述的 Apollo EM Planner 在 Frenet 坐标系下将轨迹优化分解为路径与速度两阶段，各阶段先动态规划粗解再二次规划精化，确立了解耦范式的工业基准，本文对照基线所采用的 PiecewiseJerkPathOptimizer 即源出该框架。\citet{wang2022_jfr}~报告的京东最后一公里配送车 ROVER 5.0 同样采用路径与速度独立设计再合成的解耦策略，并在多座城市完成大规模实路部署。\citet{meng2019_access}~提出格搜索与凸优化结合的解耦框架，同时明确指出解耦方法将静态与动态障碍物分置于路径与速度两个阶段处理，可能导致轨迹次优，这一论断与本文识别的结构性信息断裂相互印证。MIKU 正是针对该断裂，在不触动上述工业框架求解结构的前提下补足跨阶段信息。

\subsection{混合整数规划与端到端规划}
\label{sec:related-supplement}

除上述优化范式内的工作外，另有两类方法从不同角度处理多障碍物决策。一类是混合整数规划。\citet{kessler2023}~将避让方向等离散决策建模为整数变量，与连续轨迹优化合并为混合整数规划（Mixed-Integer Programming, MIP）在 Apollo 驾驶栈内求解；当把障碍物视为完整包络多边形时，问题进一步成为混合整数二次规划（Mixed-Integer Quadratic Programming, MIQP），可获得组合决策的全局最优，但求解复杂度为 NP 难，障碍物数量增多时实时性难以保证。MIKU 不追求组合全局最优，而是借最大间隙的局部最优性把方向选择降至 $O(k\log k)$，以可控的次优换取生产链路内的实时性。另一类是端到端学习型规划，基于深度学习与强化学习实现感知--决策--控制一体化~\citep{li2024decision,zhang2021autonomous,alahi2016social,shalevshwartz2017rss}，在场景泛化方面有所进展，但可解释性较弱、形式化安全验证困难，目前在工业级 L4 系统中多作为传统优化方法的补充。MIKU 保持优化范式，约束以显式且可检验的形式传递，并在 Apollo 11.0 生产级代码上实现，便于安全论证与工程落地。

各路线的定位对比汇总于表~\ref{tab:related}。

\begin{table}[t]
\centering
\footnotesize
\setlength{\tabcolsep}{4pt}
\caption{相关规划方法与 MIKU 的定位对比}
\label{tab:related}
\begin{tabular}{@{}lcp{3.0cm}p{2.6cm}p{2.8cm}@{}}
\toprule
代表路线 & 保留解耦 & 信息耦合方式 & 复杂度或代价 & 可行性 \\
\midrule
时空联合规划 & 否 & 三维时空直接求解 & 高，随障碍物增长 & 联合可行 \\
软代价注入 & 是 & 预测概率软代价隐式传递 & 低 & 无保证 \\
枚举--博弈 & 是 & 枚举组合后合作博弈选优 & $O(mn)$ & 选优非证明 \\
混合整数规划 & 部分 & 整数变量联合优化 & NP 难 & 组合全局最优 \\
端到端学习 & 否 & 感知--决策--控制一体化 & 训练与数据开销 & 缺形式化保证 \\
工业解耦 & 是 & 各阶段独立约束 & 低 & 仅阶段内可行 \\
MIKU（本文） & 是 & 接口处显式硬约束通道 & $O(k\log k)$ & 走廊非空时恢复 \\
\bottomrule
\end{tabular}
\end{table}

综上，时空联合规划以替换架构换取时空一致性，代价是计算开销翻倍；框架内改进或依赖软代价的隐式传递而无可行性保证，或依赖组合枚举而复杂度与改动量偏高；混合整数规划获得全局最优却受 NP 难复杂度制约；端到端学习则在优化范式之外，形式化安全验证仍是难点。MIKU 立足于工业界主流的解耦架构，在路径阶段与速度阶段的接口处建立显式硬约束传递通道，以接近不增的计算代价恢复时空协调能力，并将可行性保证收敛到可被证明的条件范围之内。

% ====== _parts/conclusion.tex ======
% ============================================================
\section{结论}
\label{sec:conclusion}
% ============================================================

本文针对路径--速度解耦架构在多障碍物动态场景下的结构性信息断裂问题，提出多障碍物交互运动学统一规划算法 MIKU（Multi-obstacle Interaction-density Kinematic-prediction Unified planner）。解耦架构是量产自动驾驶规划的主流范式，但其路径阶段不含时间维度，动态障碍物经 \texttt{PathBoundsDecider} 的静态性过滤被排除在路径决策之外，仅作为障碍物出现在速度阶段的 ST 图中，速度阶段拿到固定路径后只能被动制动。联合规划虽能恢复时空一致性，却以计算开销近乎翻倍为代价；框架内的软代价改进则使传递的信息隐式且缺乏可行性保证。MIKU 不重写求解结构，而是在解耦的两阶段接口处建立显式约束传递通道，把路径阶段对动态障碍物的时空决策以硬约束形式注入速度规划，从而在保留解耦计算优势的前提下恢复时空协调能力。

该通道由三个相互衔接的组件构成，分别回答约束传哪些、空间怎么传与时间怎么传三个问题。其一，在 SL 层以融合碰撞时间、横向重叠率、相对速度、障碍物类型与交互密度的多因子威胁度模型评估各障碍物，映射为差异化安全裕度 $\delta_i$，取值区间 $[0.10, 0.40]$\,m，使低威胁障碍物释放其占用的横向通行空间，而高威胁端不低于 Baseline 的安全水平。其二，针对组内避让方向选择，将问题建模为最大通行带宽度优化，先以扫描线按纵向重叠关系把障碍物划分为连通分量，再在组内求解最优分割；本文证明了最优分割在按安全边界排序下的连续性引理与最大间隙策略的最优性定理，将单截面搜索由 $O(2^k)$ 候选压缩至 $k+1$ 个候选分割，精确求解复杂度为 $O(k\log k)$。其三，在 ST 层以到达时间函数 $\tau(s)$ 与障碍物预测轨迹构造时变 SL 投影，进而生成时空可通行走廊，并将走廊的时间有效性约束经 \texttt{STDrivableBoundary} 注入 \texttt{PiecewiseJerkSpeedOptimizer}，与 \texttt{SpeedBoundsDecider} 的 ST 边界取交集后收紧速度上界。其中 $\tau(s)$ 由上周期速度曲线或二阶运动学模型给定，是路径阶段构建时变边界的已知输入而非优化变量，路径层的搜索空间仅为方向组合。MIKU 的整体复杂度为扫描线分组的 $O(n\log n)$ 叠加逐离散点边界构建的 $O(K)$，合计 $O(n\log n + K)$，其中 $K$ 为路径或时间离散化规模。理论层面，本文给出纯解耦规划器输出 BLOCKED 是不可行的充分而非必要条件这一命题，并证明当时空可通行走廊非空时约束增强规划器的可行性恢复保证。工程上，该通道以增量方式接入 Apollo 11.0 生产级 C++ 代码，通过 \texttt{LaneFollowPath} 新增路径边界分支并复用既有 QP 求解链路，下游控制接口无需改动。

在 Python 复现的 8 个仿真场景上的对比表明，4 个压力场景中 Baseline 与 MIKU 的通过率分别为 0/4 与 4/4，汇总 8 场景为 4/8 与 8/8；汇总平均速度由 Baseline 的 3.53\,m/s 提升至 MIKU 的 4.57\,m/s，提升 29.3\%，该提升主要源于压力场景中 Baseline 进入 \texttt{FallbackPath} 停车降级拉低了均值，并非常规巡航工况下的同等提升。在 4 个轻量微调可比场景中两者均完成通行，于同等可达前提下，MIKU 将 QP 求解总耗时由 19.16\,ms 降至 6.26\,ms，并降低了动态场景的纵向加减速峰值；压力场景中 Baseline 因降级而提前中止 QP 计算，其耗时记录不构成效率比较依据。6 变体消融显示，完整算法 M5 综合评分为基准 100.00，去除走廊注入组件的 M4 为 99.18，二者近似持平，其原因是当前压力场景未激活走廊约束的活跃边界，并非否定该组件的架构价值；其余去除关键组件的变体评分均降至 38 以下。100 次蒙特卡洛、$\pm20\%$ 权重扰动下，威胁度的最大相对偏差为 6.4\%，单因子越界概率为 0，最坏翻转率 57\% 集中于场景四中威胁度极接近的同类交通锥之间，对通行决策无实质影响。Dreamview 在环运行展示了 C++ 接入后的规划与避让行为，与上述数值指标分属两类独立证据。

\subsection{研究局限}

本文工作仍存在若干局限。验证限于 Python 数值复现与 Dreamview 在环阶段，未在真实车辆上评估，实车可用性有待在完整的感知、预测与控制链路中检验。耗时等指标取自桌面级处理器平台，车载嵌入式计算平台的实时表现需另行评估。压力场景按失效复现设计，结论应理解为特定密集与对称构型下的范围限定，而非对 Baseline 的普适优势；可比场景说明常规情形下两者的轨迹质量与速度处于同一量级。实验为确定性单次运行，尚缺大样本统计显著性检验；评价指标也未纳入最小间距等安全性度量。消融中走廊注入组件在当前压力场景未被激活，其定量贡献有待能触发活跃时间窗约束的场景补充。威胁度权重面向城市道路 30--50\,km/h 工况标定，高速大速度差场景下需重新评估。此外，$\tau(s)$ 依赖速度规划结果而速度规划又以其为输入，当前借助 Apollo 周期性重规划做单步修正，其迭代收敛性仍是开放问题。

\subsection{未来工作}

后续研究可沿四个方向展开。第一，将单向的约束传递扩展为双向反馈，把速度规划结果与自车候选行为回传至路径阶段乃至周车预测过程，把单次前向的到达时间估计推进为路径与速度的迭代求解，并对迭代收敛性与单帧计算预算给出明确界。第二，在大规模随机生成场景上开展统计验证，引入显著性检验，并量化 MIKU 相对联合规划的次优性上界。第三，引入最小间距等安全性指标，系统评估通行带收窄对碰撞裕度的影响，并给出可复现的降级与失效安全策略。第四，在多种计算平台尤其是车载嵌入式平台上评估感知噪声、预测误差与规划--控制延迟对时变通行带的影响，检验该通道在端到端时序约束下的鲁棒性。


\printbibliography[heading=bibintoc,title={参考文献}]

\end{document}
